\chapter{Energy methods}
\label{vol03:ch05}

\epigraph{No figures will be found in this work. The methods I expound
require neither constructions, nor geometrical or mechanical
arguments, but solely algebraic operations, subject to a regular and
uniform procedure.}{Joseph-Louis Lagrange, preface to \emph{M\'ecanique
Analytique} (1788), in the Boissonnade and Vagliente translation
\cite[p.~7]{hist:lagrange1788}.}

\chapmeta{Half-life: durable. Archetypes: balance, optimisation. Exercise target: 30. Project track: derivation plus simulation (hybrid). Hazard class: none. Reader path default: standard, with sections explicitly tagged. Named cases: Tacoma Narrows Bridge (1940).}

\noindent
Chapter 4 of this volume treated dynamics in the Newtonian form: a
free-body diagram for every body, a vector equation $\vec{F} = m\vec{a}$
for the translation of every centre of mass, and a moment equation for
every rotation. The method is correct and general. It is also,
for systems with many bodies, many constraints, or constraints of
complicated geometry, computationally heavy: writing the equations is
mechanical but error-prone, and the unknown internal forces between
connected bodies enter equations that the engineer then has to
eliminate. The energy-method reformulation of mechanics replaces vector
balances with scalar quantities and turns the equations of motion into
consequences of a single variational principle. The reformulation is
the work of Lagrange (1788) and Hamilton (1834); its mature engineering
form is the working tool of theoretical mechanics, structural
dynamics, robotics, and every field in which mechanical systems with
many degrees of freedom need to be analysed without bookkeeping
error. The chapter builds the scalar apparatus, works it on the
pendulum, the double pendulum, and the bead on a rotating hoop, and
closes on the failure mode the method makes possible: an analysis whose
generalised coordinates are too few to represent the motion that
actually destroys the system, the mode that brought down the Tacoma
Narrows bridge.

\section{Energy as a unifying currency}\pathtag{core}

Energy in mechanics is the integral of force through displacement,
established in chapter 4 as the work done by a force on a particle.
For a conservative force the integral depends only on the endpoints,
and the work is the change in a scalar potential energy associated
with the force. Kinetic energy is a scalar quadratic form in the
velocity. Conservation of mechanical energy in a system with only
conservative forces is the statement that the sum of kinetic and
potential energy is constant in time.

The reformulation moves three quantities to centre stage.

\engterm{Kinetic energy} $T$. For a system of particles,
\[
T = \sum_i \tfrac{1}{2} m_i \vec{v}_i \cdot \vec{v}_i,
\]
a scalar function of the velocities. For a rigid body in plane motion,
$T = \tfrac{1}{2} m v_G^2 + \tfrac{1}{2} I_G \omega^2$, as derived in
chapter 4. In all cases $T$ is a quadratic, positive-semi-definite
function of the velocities, with $T = 0$ if and only if every particle
is at rest.

\engterm{Potential energy} $U$. For conservative forces, $U$ is a
scalar function of position such that $\vec{F}_i = -\nabla_i U$, with
$\nabla_i$ the gradient with respect to particle $i$'s position. The
gravitational potential near Earth's surface is $U = mgz$ per particle;
the spring potential is $U = \tfrac{1}{2}k x^2$ per spring;
electrostatic, magnetic-dipole, and many central-force potentials have
their own standard forms.

\engterm{Generalised force}. For non-conservative forces (friction,
external applied forces, time-dependent driving forces), the system
gains a quantity that has the dimensions of force times displacement
per unit virtual displacement of a generalised coordinate; it is the
generalised force, and we will define it precisely in section 5.3
below. The non-conservative forces are then read into the formalism
without disturbing the conservative part.

The unification is not yet visible in this list. Each of the three
quantities is a scalar; combinations of them generate vector force
equations after differentiation. The variational principle of section
5.5 will state that, among all conceivable trajectories from a given
initial state to a given final state at given times, the actual
trajectory is the one that makes the time-integral of $L = T - U$
stationary. The variational form is the unifying frame; the
Newton-equivalent equations of motion emerge from it without ever
writing a vector force balance.

The engineering payoff of the reformulation is twofold. First, the
working coordinates are not Cartesian particle positions but
\emph{generalised coordinates}: any set of independent quantities
that locate the system. A pendulum has one generalised coordinate
(the swing angle), not two Cartesian particle positions; a double
pendulum has two; a robot arm has as many as its joints. Constraints
that would appear as forces in the Newtonian formulation are absorbed
into the choice of coordinates and produce no force terms in the
Lagrangian. Second, the analysis is local in the sense that bookkeeping
becomes mechanical: write $T$, write $U$, take derivatives in a
standard pattern, and the equations of motion emerge. The
error-proneness of free-body bookkeeping is replaced by the
arithmetic of partial derivatives, which is more reliable.

\begin{principle}[Mechanics in scalar form]
A mechanical system's motion can be derived from two scalar
functions of the system's state, the kinetic energy $T$ and the
potential energy $U$, together with a specification of any
non-conservative generalised forces. The equations of motion are
consequences of a single variational principle on the action
integral. No vector force balance need be drawn, and no internal
constraint force need be solved for, provided the generalised
coordinates are chosen to incorporate the constraints.
\end{principle}

A picture sharpens the claim that scalars carry the dynamics. Take the
simplest conservative system, a small-amplitude pendulum, and watch its
two energies over one period. At the turning points the bob is
momentarily at rest and all its energy sits in the gravitational
potential; at the bottom of the swing the height is least and the speed
greatest, so the energy has moved entirely into the kinetic term. The
two quantities trade back and forth, twice per period, and their sum
holds flat because gravity is conservative and the pivot does no work.
Figure~\ref{fig:vol03ch05:energy-exchange} draws the exchange. The
energy method is the discipline of reading the dynamics off this kind
of scalar accounting rather than off a force diagram.

% Figure: kinetic and potential energy exchange for the simple pendulum
% over one swing. The total stays constant; T and U trade places.
\begin{figure}[htbp]
\centering
\begin{tikzpicture}[
  axis/.style={draw=engblack, -{Stealth}, thick},
  kinetic/.style={draw=engblue, very thick},
  potential/.style={draw=engred, very thick, dashed},
  total/.style={draw=englime, very thick},
  guide/.style={draw=enggray, dashed, thin},
  label/.style={font=\small},
]
% x = phase from 0 to 2pi (plotted 0..10), y = energy fraction 0..1 (plotted 0..5)
\draw[axis] (0,0) -- (10.6,0) node[anchor=west, label] {phase $\omega t$};
\draw[axis] (0,0) -- (0,5.6)  node[anchor=south, label] {energy / $E$};

% x ticks: 0, pi/2, pi, 3pi/2, 2pi (tick marks)
\foreach \x in {0, 2.5, 5, 7.5, 10} {
  \draw (\x,-0.08) -- (\x,0.08);
}
% x tick labels placed individually to keep slashes out of the foreach parser
\node[font=\scriptsize, below] at (0,-0.12) {$0$};
\node[font=\scriptsize, below] at (2.5,-0.12) {$\pi/2$};
\node[font=\scriptsize, below] at (5,-0.12) {$\pi$};
\node[font=\scriptsize, below] at (7.5,-0.12) {$3\pi/2$};
\node[font=\scriptsize, below] at (10,-0.12) {$2\pi$};
% y ticks: 0, 0.5, 1.0
\foreach \y/\h in {0/0, 2.5/0.5, 5/1.0} {
  \draw (-0.08,\y) -- (0.08,\y);
  \node[font=\scriptsize, left] at (-0.12,\y) {$\h$};
}

% Total energy: constant at E (y = 5)
\draw[total] (0,5) -- (10,5);
\node[englime, font=\scriptsize, anchor=west] at (10.05, 5) {$T+U$};

% Kinetic energy: T/E = sin^2(omega t / 2) ... use small-amplitude:
% theta = theta0 cos(omega t), thetadot = -theta0 omega sin(omega t)
% T = (1/2) m L^2 thetadot^2 propto sin^2(omega t); U propto cos^2(omega t)
% phase variable s = omega t. x = 5*s/pi, so s = x*pi/5.
\draw[kinetic] plot[domain=0:10, samples=200]
  ({\x}, {5 * (sin(deg(\x*3.14159/5)))^2});
\node[engblue, font=\scriptsize, anchor=south] at (2.5, 5.15) {$T$};

% Potential energy: U/E = cos^2(omega t)
\draw[potential] plot[domain=0:10, samples=200]
  ({\x}, {5 * (cos(deg(\x*3.14159/5)))^2});
\node[engred, font=\scriptsize, anchor=north] at (5.0, 4.7) {$U$};

\end{tikzpicture}
\caption[Kinetic and potential energy exchange in the simple
pendulum]{Energy exchange in a small-amplitude simple pendulum over
one full period. The kinetic energy $T$ (blue) is maximal at the
bottom of the swing where the speed is greatest; the potential energy
$U$ (red dashed) is maximal at the turning points where the bob is
instantaneously at rest. The two trade places twice per period, and
their sum, the total mechanical energy $E = T+U$ (green), stays
constant because the only force doing work is conservative gravity and
the pivot constraint is workless. The energy method takes this scalar
bookkeeping as its starting point: write $T$, write $U$, and let the
equations of motion follow from how the two are differentiated.}
\label{fig:vol03ch05:energy-exchange}
\end{figure}


A second picture fixes the relation between a potential and the force
it encodes. For a single coordinate the conservative force is the
negative slope of the potential, $F = -dU/dx$, so a potential drawn as
a curve carries the whole force law in its shape: equilibria are its
stationary points, the restoring force at any displacement is its local
slope, and the stiffness is its curvature. The spring potential
$U = \tfrac12 k x^2$ is the canonical well, drawn in
figure~\ref{fig:vol03ch05:spring-potential}; a stiffer spring is a
steeper well, and the bottom of the well is the stable equilibrium to
which the mass returns. Most of the analysis in this chapter is the
multi-coordinate generalisation of reading a curve: locate the
stationary points of an energy, examine the curvature there, and let
the equilibria and their stability follow.

% Figure: potential-energy well of a linear spring and the force as its
% negative gradient. The parabola is U(x) = (1/2) k x^2; the slope at a
% point is the restoring force.
\begin{figure}[htbp]
\centering
\begin{tikzpicture}[
  axis/.style={draw=engblack, -{Stealth}, thick},
  well/.style={draw=engblue, very thick},
  force/.style={draw=engred, -{Stealth}, very thick},
  guide/.style={draw=enggray, dashed, thin},
  label/.style={font=\small},
]
% x = displacement from -4 to 4, y = U from 0 to 5
\draw[axis] (-4.4,0) -- (4.4,0) node[anchor=west, label] {$x$};
\draw[axis] (0,-0.3) -- (0,5.4) node[anchor=south, label] {$U(x)$};

\foreach \x in {-3,-2,-1,1,2,3} {
  \draw (\x,-0.07) -- (\x,0.07);
}
\node[font=\scriptsize, below] at (3,-0.1) {$x>0$};
\node[font=\scriptsize, below] at (-3,-0.1) {$x<0$};

% Parabola U = 0.31 x^2 (scaled so U(4)=5)
\draw[well] plot[domain=-4:4, samples=120] ({\x}, {0.3125*\x*\x});
\node[engblue, label, anchor=south west] at (2.6, 2.6) {$U=\tfrac12 k x^2$};

% A point on the right side at x = 2.4, U = 1.8
% slope dU/dx = 0.625 x = 1.5; force F = -dU/dx points toward origin (left)
\fill[engblack] (2.4, 1.8) circle (2pt);
% tangent line at the point for slope illustration
\draw[guide] (1.4, {1.8 - 1.5*1.0}) -- (3.4, {1.8 + 1.5*1.0});
% restoring force arrow at the bottom, pointing left (toward equilibrium)
\draw[force] (2.4, 0.45) -- (1.2, 0.45);
\node[engred, font=\scriptsize, anchor=south] at (1.8, 0.5) {$F=-\dfrac{dU}{dx}$};

% Equilibrium marker
\fill[englime] (0,0) circle (2.2pt);
\node[englime, font=\scriptsize, anchor=north east] at (-0.1,-0.05) {equilibrium};

\end{tikzpicture}
\caption[Spring potential well and the restoring force]{The potential
energy of a linear spring, $U(x)=\tfrac12 k x^2$, drawn as a parabolic
well (blue). The restoring force is the negative gradient,
$F=-dU/dx=-kx$, so at any displaced point (black dot) the force (red
arrow) points back toward the minimum at $x=0$ (green). The steeper the
well, the larger the stiffness $k$ and the stronger the restoring
force for a given displacement. Reading mechanics off a potential is
the working habit the energy method installs: equilibria sit at
stationary points of $U$, stability is the curvature of $U$ there, and
the force at any point is the local slope.}
\label{fig:vol03ch05:spring-potential}
\end{figure}


\begin{estimation}
\textbf{Question.} A child on a $2.5\,\si{\meter}$ playground swing
starts from rest with the ropes horizontal (a $90^{\circ}$ amplitude)
and is brought to rest by a parent at the bottom of the first descent.
Estimate the impulse the parent must supply and the average force over a
half-second catch, and check whether energy methods or a force balance
is the shorter route.

\textbf{Working.} Take the child plus seat as a point mass of
$m \approx 25\,\si{\kilogram}$. The drop from horizontal ropes to the
bottom is the rope length, $h = L = 2.5\,\si{\meter}$. Energy
conservation gives the speed at the bottom without any path integration:
$\tfrac12 m v^2 = m g h$, so
\[
v = \sqrt{2 g h} = \sqrt{2 \cdot 9.81 \cdot 2.5}
\approx 7.0\,\si{\meter\per\second}.
\]
The momentum to be removed is $p = m v \approx 25 \cdot 7.0
= 175\,\si{\kilogram\meter\per\second}$, which is the impulse required.
Over a catch lasting $\Delta t \approx 0.5\,\si{\second}$ the average
force is
\[
F_\text{avg} \approx \frac{p}{\Delta t}
= \frac{175}{0.5} = 350\,\si{\newton},
\]
roughly the child's weight again ($mg \approx 245\,\si{\newton}$)
plus enough to decelerate, a few hundred newtons, demanding but
manageable for an adult. A force-balance route would have required
integrating the tension and the tangential weight component along the
quarter-circle arc to reach the same speed; energy conservation read it
off the height drop in one line. The number to carry away is the order
of magnitude: catching a swing at full amplitude is a several-hundred-
newton job, and the energy method is the faster path to it.
\end{estimation}

\section{Conservation of energy in mechanics}\pathtag{standard}

For a system whose forces are all conservative and whose constraints
do no work (a constraint is workless if its reaction force is
perpendicular to the permitted virtual displacement, which is true of
smooth pin joints, smooth surfaces, and inextensible ropes), the
total mechanical energy is constant in time:
\[
E = T + U = \text{constant}.
\]
The result follows from Newton's second law $\vec{F} = m\vec{a}$ and
the identity $\vec{a}\cdot d\vec{r} = d(\tfrac{1}{2} v^2)$, as derived
in chapter 4. For a system that includes non-conservative forces
(friction, drag, applied driving), energy is not conserved; the rate
of change of $T + U$ equals the rate at which non-conservative forces
do work on the system.

\begin{archetype}[Balance]
Conservation of energy is the mechanical instance of the balance
archetype: a quantity is neither created nor destroyed, so its
accumulation in a system equals the net flow across the boundary. The
same accounting governs conservation of mass in a control volume
(Volume IV), conservation of charge at a circuit node (Kirchhoff's
current law), conservation of momentum in a collision, and the first
law of thermodynamics, where the change in internal energy balances heat
added and work done. The recurring discipline is to draw a boundary,
list every inflow and outflow of the conserved quantity, and set the
accumulation equal to the net flux. In a conservative mechanical system
the boundary flux is zero and the balance collapses to $T + U =
\text{constant}$; the energy method generalises this from a constant to
a rate equation when non-conservative work crosses the boundary.
\end{archetype}

The principle of \engterm{virtual work} is the static-equilibrium
limit of energy methods. For a system in static equilibrium, any
virtual displacement consistent with the constraints (a hypothetical
infinitesimal motion that respects the geometric constraints but is
not necessarily a real motion of the system) produces zero virtual
work by the applied forces:
\[
\delta W = \sum_i \vec{F}_i \cdot \delta \vec{r}_i = 0.
\]
The sum is over the applied forces, not over the constraint forces,
because workless constraints produce no virtual work by definition.
The principle is equivalent to the statement that, in static
equilibrium with workless constraints, the applied forces' resultant
in every allowed direction of motion is zero.

For conservative applied forces, virtual work has the simpler form
\[
\delta U = 0,
\]
the variation of the potential energy under any allowed virtual
displacement vanishes. The applied forces' equilibrium is the
stationarity of $U$ under variations consistent with the constraints.
The result was established by Lagrange in the 1788 \emph{M\'ecanique
Analytique} as the generalisation of the equilibrium conditions to
arbitrary mechanical systems \cite{hist:lagrange1788}.

\textbf{Worked example: the lever by virtual work.} A rigid lever
pivots without friction at a fulcrum, with an applied force $F_1$ at
distance $a$ on one side and a resisting force $F_2$ at distance $b$ on
the other (figure~\ref{fig:vol03ch05:virtual-work}). The Newtonian
route writes a moment balance about the fulcrum and, if the full state
is wanted, a force balance that introduces the unknown pin reaction.
The virtual-work route skips the reaction. Give the lever a virtual
rotation $\delta\phi$ about the fulcrum, the one motion the pin
constraint permits. The two ends move through $\delta r_1 = a\,\delta\phi$
and $\delta r_2 = b\,\delta\phi$. The pin reaction does no virtual work
because the fulcrum point does not move. Setting the total virtual work
of the applied forces to zero,
\[
\delta W = F_1\,a\,\delta\phi - F_2\,b\,\delta\phi = 0
\quad\Longrightarrow\quad
F_1 a = F_2 b,
\]
the law of the lever, with the reaction never computed. The economy is
the same one the Lagrangian formulation delivers in dynamics: workless
constraints drop out, and a scalar balance survives.

% Figure: virtual work on a simple lever / two-bar linkage. A virtual
% displacement of the input rotates the mechanism; the principle of
% virtual work balances input and output through the geometry alone, with
% no need to compute the pin reaction.
\begin{figure}[htbp]
\centering
\begin{tikzpicture}[
  bar/.style={draw=engblack, very thick},
  ghost/.style={draw=enggray, very thick, dashed},
  force/.style={draw=engred, -{Stealth}, very thick},
  pin/.style={draw=engblack, fill=white, circle, inner sep=1.6pt},
  vd/.style={draw=engblue, -{Stealth}, thick},
  label/.style={font=\small},
]
% Fulcrum at origin
\node[pin] (F) at (0,0) {};
\node[font=\scriptsize, below] at (0,-0.15) {fulcrum};

% Lever: left arm length a = 3.4 (input), right arm length b = 1.7 (output)
% horizontal at rest
\coordinate (A) at (-3.4, 0);
\coordinate (B) at (1.7, 0);
\draw[bar] (A) -- (B);

% Virtual rotation delta-phi: ends move vertically by a*dphi and b*dphi
% rotate by small angle (drawn at 12 deg for visibility)
\coordinate (Ap) at ({-3.4*cos(12)}, {-3.4*sin(12)});
\coordinate (Bp) at ({1.7*cos(12)}, {1.7*sin(12)});
\draw[ghost] (Ap) -- (Bp);

% input force F1 downward at A (load pushes A down through delta r1 = a dphi)
\draw[force] (-3.4, 1.3) -- (-3.4, 0.15);
\node[engred, font=\small, anchor=south] at (-3.4, 1.35) {$F_1$};
% output force F2 downward at B resisting
\draw[force] (1.7, -1.3) -- (1.7, -0.15);
\node[engred, font=\small, anchor=north] at (1.7, -1.35) {$F_2$};

% virtual displacements
\draw[vd] (-3.4, 0) -- (Ap);
\node[engblue, font=\scriptsize, anchor=east] at ($(Ap)+(-0.05,0)$) {$\delta r_1 = a\,\delta\phi$};
\draw[vd] (1.7, 0) -- (Bp);
\node[engblue, font=\scriptsize, anchor=west] at ($(Bp)+(0.05,0)$) {$\delta r_2 = b\,\delta\phi$};

% arc showing delta phi at fulcrum
\draw[draw=englime, -{Stealth}, thick] (0.9,0) arc (0:12:0.9);
\node[englime, font=\scriptsize, anchor=west] at (0.95, 0.15) {$\delta\phi$};

\end{tikzpicture}
\caption[Principle of virtual work on a lever]{The principle of virtual
work applied to a rigid lever pivoted at a frictionless fulcrum. A
virtual rotation $\delta\phi$ (green) consistent with the pin constraint
moves the left end down by $\delta r_1=a\,\delta\phi$ and the right end
up by $\delta r_2=b\,\delta\phi$ (blue). The pin reaction does no virtual
work because the fulcrum does not move. Setting the total virtual work
of the applied forces to zero, $F_1\,a\,\delta\phi - F_2\,b\,\delta\phi=0$,
gives the equilibrium law $F_1 a = F_2 b$ directly, without ever solving
for the reaction at the pin. Static equilibrium is the zero-velocity
limit of the energy method: the workless constraints drop out, leaving a
scalar balance read off the geometry.}
\label{fig:vol03ch05:virtual-work}
\end{figure}


\textbf{Worked example: speed from energy conservation.} A block of
mass $m$ released from rest slides without friction down a smooth track
and drops a height $h$. No force other than gravity does work (the
normal force is perpendicular to the motion and so is workless), so
$T + U$ is constant. Taking the release point as the height datum,
$0 + mgh = \tfrac12 m v^2 + 0$ at the bottom, hence
\[
v = \sqrt{2 g h},
\]
independent of the mass and of the shape of the track. The same result
extracted by Newton's method would require resolving the weight along a
track whose slope changes continuously and integrating the tangential
acceleration over the path; energy conservation reads the answer off
the endpoints. This is the recurring trade the chapter develops: a
scalar bookkeeping of energy replaces a vector integration along a
path, and the path's geometry, which dominates the Newtonian
calculation, drops out when only the endpoints matter.

D'Alembert's principle extends virtual work to dynamics: for a system
not in static equilibrium, the applied forces minus the inertial
``forces'' $-m\vec{a}$ produce zero virtual work for any virtual
displacement consistent with the constraints:
\[
\sum_i (\vec{F}_i - m_i \vec{a}_i) \cdot \delta \vec{r}_i = 0.
\]
The inertial term $-m\vec{a}$ is sometimes called a pseudo-force or
an inertial force; its appearance in the equation captures Newton's
second law's content as a special case of virtual work. The
formulation is the bridge from Newtonian dynamics to Lagrangian
mechanics: substituting $T$ and $U$ for the kinematic and force terms
turns d'Alembert's principle into the Euler-Lagrange equations of
the next section.

\section{Lagrangian mechanics introduced}\pathtag{core}

A mechanical system with $n$ generalised coordinates $q_1, q_2, \ldots,
q_n$ (independent quantities that locate the system completely, given
the constraints) has a kinetic energy $T = T(q, \dot{q}, t)$ and a
potential energy $U = U(q, t)$. The \engterm{Lagrangian} is
\[
L(q, \dot{q}, t) = T(q, \dot{q}, t) - U(q, t).
\]
The Lagrangian is a scalar function of $2n + 1$ arguments: the $n$
generalised coordinates, the $n$ generalised velocities, and time.

The equations of motion are the \engterm{Euler-Lagrange equations}:
for each generalised coordinate $q_k$, $k = 1, \ldots, n$,
\[
\frac{d}{dt}\left(\frac{\partial L}{\partial \dot{q}_k}\right) -
\frac{\partial L}{\partial q_k} = Q_k^{\text{nc}},
\]
where $Q_k^{\text{nc}}$ is the generalised non-conservative force
conjugate to $q_k$ (zero if all forces are conservative). The
left-hand side is constructed from the Lagrangian in a mechanical way:
take the partial derivative of $L$ with respect to the generalised
velocity, then differentiate the result with respect to time; subtract
the partial derivative of $L$ with respect to the corresponding
generalised coordinate. The result, set equal to the generalised
non-conservative force, is the equation of motion for that coordinate.

The Euler-Lagrange equations are mathematically equivalent to
Newton's second law applied to every particle, with constraint forces
solved for. The equivalence is provable: starting from d'Alembert's
principle and a coordinate transformation from Cartesian particle
positions to generalised coordinates, the inertial term reduces
algebraically to $d/dt(\partial T/\partial \dot{q}) - \partial T/
\partial q$, and the conservative force terms to $-\partial U/\partial
q$. The Lagrangian formulation is therefore not a new physics; it is
a more economical bookkeeping of the same physics.

The economy is the key engineering benefit. A robot arm with $n$
joints has $n$ generalised coordinates (the joint angles); writing
its equations of motion in Newtonian form requires drawing $n + 1$
free-body diagrams (one per link plus the base reactions) and solving
$3n + 3$ equations in the joint reaction forces, then eliminating
$2n$ unknown reaction components. The Lagrangian form writes one
scalar $L$, takes $n$ partial-derivative pairs, and arrives at $n$
equations of motion directly. For $n > 2$ the difference is large.

\engterm{Cyclic coordinates and conserved momenta}. A generalised
coordinate that does not appear in $L$ (only $\dot{q}_k$ appears) is
\engterm{cyclic}. For cyclic $q_k$, the Euler-Lagrange equation
reduces to
\[
\frac{d}{dt}\left(\frac{\partial L}{\partial \dot{q}_k}\right) = 0,
\]
which states that the corresponding generalised momentum $p_k =
\partial L / \partial \dot{q}_k$ is constant. Identifying cyclic
coordinates is the standard route to conservation laws: a system
whose $L$ does not depend on a Cartesian coordinate conserves the
corresponding linear momentum; a system whose $L$ does not depend on
a rotation angle conserves the corresponding angular momentum. The
connection between symmetries of $L$ and conservation laws is the
theorem of Emmy Noether (1918), which extends to continuum mechanics
and field theory in the standard way.

\textbf{Worked example: a bead on a frictionless wire.} The smallest
non-trivial application makes the bookkeeping concrete. A bead of mass
$m$ slides on a straight frictionless horizontal wire with no applied
force. The one generalised coordinate is the position $x$ along the
wire. The kinetic energy is $T = \tfrac12 m \dot{x}^2$, the potential
energy is constant and may be set to zero, so $L = \tfrac12 m \dot{x}^2$.
The coordinate $x$ does not appear in $L$, only $\dot{x}$ does, so $x$
is cyclic and the conjugate momentum $p = \partial L/\partial\dot{x} =
m\dot{x}$ is conserved. The Euler-Lagrange equation confirms it:
$\tfrac{d}{dt}(m\dot{x}) - 0 = 0$, hence $m\ddot{x} = 0$, the bead
moves at constant velocity. The conserved momentum is the linear
momentum, and its conservation is the symmetry of the Lagrangian under
translation along the wire. The same logic, applied to a rotation
coordinate that is absent from $L$, conserves angular momentum; applied
to time, when $L$ has no explicit time dependence, it conserves energy.
Conservation laws are read off the absences in the Lagrangian.

\begin{mastery}
\textbf{The Hamiltonian and phase space.} The Lagrangian formulation
works in the $2n$-dimensional space of coordinates and velocities
$(q, \dot{q})$. The Hamiltonian formulation, due to Hamilton (1834) and
developed by Jacobi, trades each velocity $\dot{q}_k$ for its conjugate
momentum $p_k = \partial L/\partial\dot{q}_k$ and works in the phase
space of $(q, p)$. The exchange is a Legendre transform, the same
construction that carries internal energy to enthalpy in
thermodynamics:
\[
H(q, p, t) = \sum_k p_k \dot{q}_k - L(q, \dot{q}, t),
\]
with the velocities eliminated in favour of the momenta. For the broad
class of systems whose kinetic energy is a homogeneous quadratic in the
velocities and whose potential is velocity-independent, $H = T + U$,
the total energy. The equations of motion become Hamilton's canonical
equations, a symmetric first-order pair,
\[
\dot{q}_k = \frac{\partial H}{\partial p_k},
\qquad
\dot{p}_k = -\frac{\partial H}{\partial q_k}.
\]
The pair has $2n$ first-order equations in place of $n$ second-order
ones, a form that suits numerical integration, that exposes the
geometry of phase space (Liouville's theorem on the conservation of
phase-space volume, the symplectic structure preserved by the flow),
and that is the bridge to statistical mechanics and to the canonical
quantisation of quantum mechanics. The Landau-Lifshitz treatment
develops the canonical formalism from the action principle with minimal
apparatus \cite{text:landau-mechanics}. For the working analysis in
this chapter the Lagrangian form is sufficient; the Hamiltonian form is
the entry point to the deeper structure and to the symplectic
integrators used when long-time energy conservation matters.
\end{mastery}

\section{Generalised coordinates and constraints}\pathtag{standard}

A mechanical system in three-dimensional space with $N$ particles
has $3N$ Cartesian coordinates. Constraints reduce the system's
freedom. A \engterm{holonomic} constraint is an equation of the form
\[
f(r_1, r_2, \ldots, r_N, t) = 0,
\]
relating the particle positions and time. Each independent holonomic
constraint reduces the degrees of freedom by one. A system with $N$
particles and $c$ holonomic constraints has $n = 3N - c$ degrees of
freedom, and $n$ independent generalised coordinates suffice to
locate it.

Examples of holonomic constraints. A particle confined to a surface
$f(x, y, z) = 0$ has one constraint, two degrees of freedom; the
surface coordinates are the natural generalised coordinates. A bead
on a rotating circular hoop has the constraint $x^2 + y^2 = R^2$
combined with the rotating-frame angle; one generalised coordinate
(the bead's angular position on the hoop) suffices. A double pendulum
in two dimensions has four Cartesian coordinates (two per bob) and
two constraints (each bob's distance from its pivot is constant);
two generalised coordinates (the two swing angles) suffice.

A \engterm{nonholonomic} constraint involves the velocities and
cannot be integrated into a relation among the coordinates alone. The
canonical example is a wheel rolling without slipping on a surface:
the rolling condition $\vec{v}_{\text{contact}} = 0$ relates the
wheel's velocity to its angular velocity, but the integrated form is
a path-dependent relation that does not reduce to a constraint on the
coordinates. Nonholonomic systems retain more generalised
coordinates than their kinematic degrees of freedom suggest; the
analysis uses Lagrange multipliers or the Gibbs-Appell formulation,
both beyond this chapter's scope. The robotics and vehicle-dynamics
treatments in Volume VIII (chapter 8) develop the nonholonomic case.

The choice of generalised coordinates is the engineer's. Independent
coordinates that completely locate the system are required; beyond
that, the choice is tactical. A pendulum's swing angle $\theta$ is
the natural coordinate; the bob's Cartesian position is allowed but
introduces a constraint $x^2 + y^2 = L^2$ that the swing angle
absorbs automatically. A system whose constraints are mismatched
with its coordinate choice produces algebraic complication; matched
coordinates produce algebraic simplicity. The standard heuristic is
that coordinates should be chosen so that each constraint is
identically satisfied by the chosen coordinate set.

\engterm{Degrees of freedom}, the count of independent generalised
coordinates required, is the most important number to determine
correctly before any energy-method analysis. A wrong count produces
the failure mode of section 5.7: an analysis whose equations of
motion miss a mode of the real system.

\section{Hamilton's principle}\pathtag{standard}

Among all conceivable trajectories $q(t)$ from an initial state $q
(t_1) = q_1$ to a final state $q(t_2) = q_2$ at fixed times $t_1$
and $t_2$, the trajectory the system actually follows is the one
that makes the \engterm{action}
\[
S = \int_{t_1}^{t_2} L(q, \dot{q}, t) \,dt
\]
stationary under arbitrary variations of $q(t)$ with the endpoints
fixed. This is \engterm{Hamilton's principle}, established by William
Rowan Hamilton in his 1834 paper ``On a General Method in Dynamics.''
Figure~\ref{fig:vol03ch05:action-paths} draws the construction: among
all admissible trial paths sharing the two endpoints, the physical path
is the one whose action does not change to first order under a small
deformation.

% Figure: Hamilton's principle. Among all trial paths between fixed
% endpoints, the physical path makes the action stationary. Trial paths
% are drawn as perturbations of the true path.
\begin{figure}[htbp]
\centering
\begin{tikzpicture}[
  axis/.style={draw=engblack, -{Stealth}, thick},
  truepath/.style={draw=engblue, very thick},
  trial/.style={draw=enggray, thick, dashed},
  endpt/.style={draw=engred, fill=engred, circle, inner sep=1.8pt},
  label/.style={font=\small},
]
% t axis horizontal, q axis vertical
\draw[axis] (0,0) -- (10.4,0) node[anchor=west, label] {$t$};
\draw[axis] (0,0) -- (0,5.4) node[anchor=south, label] {$q$};

% endpoints fixed: (t1, q1) = (1, 1) plotted (1,1); (t2, q2) = (9, 4)
\node[endpt] (P1) at (1,1) {};
\node[endpt] (P2) at (9,4) {};
\node[font=\scriptsize, anchor=north east] at (1,1) {$(t_1,q_1)$};
\node[font=\scriptsize, anchor=north west] at (9,4) {$(t_2,q_2)$};
\draw[draw=enggray, dashed, thin] (1,1) -- (1,0) node[font=\scriptsize, below] {$t_1$};
\draw[draw=enggray, dashed, thin] (9,4) -- (9,0) node[font=\scriptsize, below] {$t_2$};

% true path: smooth curve from (1,1) to (9,4); use a gentle parabola-ish arc
\draw[truepath] plot[domain=1:9, samples=80]
  ({\x}, {1 + 0.375*(\x-1) + 0.6*sin(deg((\x-1)*3.14159/8))});
\node[engblue, label, anchor=south] at (5, 4.4) {physical path (action stationary)};

% trial path 1: above
\draw[trial] plot[domain=1:9, samples=80]
  ({\x}, {1 + 0.375*(\x-1) + 0.6*sin(deg((\x-1)*3.14159/8)) + 0.9*sin(deg((\x-1)*3.14159/8))});
% trial path 2: below
\draw[trial] plot[domain=1:9, samples=80]
  ({\x}, {1 + 0.375*(\x-1) + 0.6*sin(deg((\x-1)*3.14159/8)) - 0.9*sin(deg((\x-1)*3.14159/8))});
\node[font=\scriptsize, enggray, anchor=west] at (5.2, 2.0) {trial paths $q+\delta q$};

% variation arrow at t = 5
\draw[draw=englime, {Stealth}-{Stealth}, thick] (5, {1+0.375*4+0.6+(-0.9)}) -- (5, {1+0.375*4+0.6+0.9});
\node[englime, font=\scriptsize, anchor=west] at (5.05, {1+0.375*4+0.6}) {$\delta q$};

\end{tikzpicture}
\caption[Hamilton's principle and admissible trial paths]{Hamilton's
principle in one coordinate. The endpoints $(t_1,q_1)$ and $(t_2,q_2)$
are fixed (red). Among all admissible trial paths $q(t)+\delta q(t)$
that share those endpoints (grey dashed), the path the system actually
follows (blue) is the one for which the action $S=\int_{t_1}^{t_2}L\,dt$
is stationary: $\delta S=0$ to first order in the variation $\delta q$
(green). The Euler-Lagrange equations are exactly the condition that
this stationarity hold for every admissible variation. Mechanics, cast
this way, is an optimisation problem, and the same variational frame
carries into elasticity, optics, and field theory.}
\label{fig:vol03ch05:action-paths}
\end{figure}


The Euler-Lagrange equations are the necessary conditions for $S$ to
be stationary. Varying the action with respect to a small perturbation
$\delta q(t)$ that vanishes at the endpoints,
\[
\delta S = \int_{t_1}^{t_2} \left[\frac{\partial L}{\partial q}
\delta q + \frac{\partial L}{\partial \dot{q}} \delta \dot{q}\right]
\,dt.
\]
Integrating the second term by parts (the boundary terms vanish
because $\delta q$ vanishes at the endpoints):
\[
\delta S = \int_{t_1}^{t_2} \left[\frac{\partial L}{\partial q} -
\frac{d}{dt}\frac{\partial L}{\partial \dot{q}}\right] \delta q
\,dt.
\]
For $\delta S = 0$ to hold for arbitrary $\delta q(t)$, the bracketed
expression must vanish at every $t$:
\[
\frac{d}{dt}\frac{\partial L}{\partial \dot{q}} - \frac{\partial L}{
\partial q} = 0,
\]
the Euler-Lagrange equation. The same derivation in $n$ coordinates
gives the $n$ Euler-Lagrange equations of section 5.3.

The variational form is more than a derivation route. It expresses
mechanics as an optimisation: the trajectory the system follows is
optimal among all conceivable trajectories, with respect to a
specific functional (the action integral). The optimisation
interpretation generalises beyond classical mechanics. Hamilton's
principle is the starting point for the modern formulation of
continuum mechanics, fluid dynamics, electromagnetism, general
relativity, and quantum field theory; the action principle is the
single most general way to specify a physical theory.

The engineering payoff is, again, economy. Once the system's
Lagrangian is identified, the equations of motion follow mechanically.
The variational form also gives a route to approximate solutions: an
assumed trial trajectory $q(t; a_1, a_2, \ldots, a_m)$ with $m$
adjustable parameters $a_k$ produces an action $S(a_1, a_2, \ldots,
a_m)$ that is a function (not a functional) of the $m$ parameters;
making $S$ stationary with respect to the parameters yields an
approximate equation of motion that is exact within the trial family
and is the basis of the Rayleigh-Ritz method, the finite-element
method, and most numerical solvers in structural dynamics.

\begin{mastery}
\textbf{The Rayleigh quotient as an energy estimate.} The variational
structure gives a fast, conservative estimate of a vibrating system's
lowest natural frequency without solving the equations of motion. For a
linear conservative system with a trial mode shape, the ratio of peak
potential energy to a reference kinetic energy bounds the fundamental
frequency from above. In the continuous case of a beam or string, the
Rayleigh quotient is
\[
\omega^2 \le R[\phi]
= \frac{\displaystyle\int (\text{stiffness})\,(\phi')^2\,dx}
       {\displaystyle\int (\text{mass density})\,\phi^2\,dx},
\]
evaluated on any admissible trial shape $\phi$ that satisfies the
geometric boundary conditions. The estimate is always an over-estimate
of the true fundamental, and it is remarkably accurate even for crude
trial shapes, because the quotient is stationary at the true mode and
so is insensitive to first-order error in $\phi$. A parabolic trial
shape for a cantilever, for instance, lands within a few per cent of
the exact fundamental frequency. The method is the practising
structural engineer's back-of-envelope check on a finite-element
natural-frequency result: if the code reports a fundamental far from the
Rayleigh estimate, the model, not the structure, is usually at fault.
The same quotient, generalised to a matrix eigenvalue problem,
underlies the power-iteration and subspace methods that the
finite-element solvers themselves use. Den~Hartog's vibrations text
develops the Rayleigh method as the working tool for hand estimates of
natural frequencies \cite{text:den-hartog-vibrations}.
\end{mastery}

\begin{archetype}[Optimisation]
Variational principles cast mechanics as optimisation: among
admissible trajectories, the physical one extremises a scalar
functional. The same archetype recurs in thermodynamics (equilibrium
minimises Gibbs free energy), in chemistry (reactions follow
free-energy gradients), in optics (Fermat's principle: light follows
the path of stationary optical length), in elasticity (the deformed
state minimises elastic potential energy), in finite-element
analysis (the discretised approximation minimises a discrete energy
functional), and in machine learning (training minimises a loss
function). The optimisation archetype is one of the unifying
patterns of engineering mathematics; chapter 4 of Volume IX develops
it as a stand-alone framework.
\end{archetype}

\section{Worked examples}\pathtag{standard}

\subsection{Simple pendulum}

A simple pendulum consists of a bob of mass $m$ suspended from a
fixed pivot by a massless inextensible string of length $L$. The
generalised coordinate is the swing angle $\theta$, measured from
the downward vertical.

\textbf{Kinetic energy.} The bob's velocity is $L\dot{\theta}$
tangent to its circular path, so $T = \tfrac{1}{2} m L^2 \dot{\theta}^2$.

\textbf{Potential energy.} With $U = 0$ at the lowest point of the
swing, the bob's height above the lowest point is $L(1 - \cos\theta)$,
so $U = m g L (1 - \cos\theta)$.

\textbf{Lagrangian.}
\[
L = T - U = \tfrac{1}{2} m L^2 \dot{\theta}^2 - m g L (1 - \cos\theta).
\]

\textbf{Euler-Lagrange equation.}
\[
\frac{\partial L}{\partial \dot{\theta}} = m L^2 \dot{\theta}, \quad
\frac{d}{dt}\frac{\partial L}{\partial \dot{\theta}} = m L^2
\ddot{\theta}, \quad \frac{\partial L}{\partial \theta} = -m g L
\sin\theta.
\]
Therefore $m L^2 \ddot{\theta} + m g L \sin\theta = 0$, or
\[
\ddot{\theta} + \frac{g}{L}\sin\theta = 0.
\]
This is the equation of motion of a simple pendulum, identical to
the result of a Newtonian analysis with free-body diagrams of the
bob and constraint reactions at the pivot. The Lagrangian derivation
involved one kinetic-energy scalar, one potential-energy scalar, and
mechanical differentiation; the Newtonian derivation requires the
free-body diagram, the resolution of the tension in the string into
radial and tangential components, and the explicit elimination of
the unknown tension to leave the tangential equation.

For small amplitudes ($\theta \ll 1$), $\sin\theta \approx \theta$ and
the equation linearises to $\ddot{\theta} + (g/L)\theta = 0$, the
simple harmonic oscillator with period $T = 2\pi\sqrt{L/g}$. A
$1\,\si{\meter}$ pendulum has period approximately $2.0\,\si{\second}$.

At larger amplitude the small-angle period under-predicts the truth,
because the restoring torque $-mgL\sin\theta$ grows more slowly than the
linear law $-mgL\,\theta$ that the small-angle formula assumes. The
leading correction is $T \approx T_0\,(1 + \theta_0^2/16)$, and the
exact period is the complete elliptic integral of the first kind. The
script \path{code/pendulum_period.py} integrates the full nonlinear
equation, measures the period by detecting the return to the turning
point, and writes the comparison to \path{data/pendulum_period.csv}: at
$\theta_0 = 45^\circ$ the small-angle formula is low by about
$4\,\%$ and the first correction is good to better than $0.2\,\%$;
at $\theta_0 = 90^\circ$ the small-angle formula is low by $18\,\%$
and even the first correction is short by $2\,\%$. The exercises return
to these numbers.

\subsection{Double pendulum}

A double pendulum has two bobs of masses $m_1$ and $m_2$ on
inextensible massless strings of lengths $L_1$ and $L_2$, the second
suspended from the first. Generalised coordinates: the swing angles
$\theta_1$ (upper) and $\theta_2$ (lower), both measured from the
downward vertical (figure~\ref{fig:vol03ch05:double-pendulum}).

% Figure: double-pendulum geometry and its two generalised coordinates.
% Two massless links, two bobs, angles measured from the downward
% vertical at each pivot.
\begin{figure}[htbp]
\centering
\begin{tikzpicture}[
  link/.style={draw=engblack, very thick},
  bob/.style={draw=engblue, fill=engblue, circle, inner sep=2.5pt},
  vert/.style={draw=enggray, dashed, thin},
  ang/.style={draw=engred, -{Stealth}, thick},
  label/.style={font=\small},
]
% Pivot at origin (top)
\fill[engblack] (0,0) circle (1.6pt);
\node[font=\scriptsize, anchor=south] at (0,0.1) {pivot};

% Upper link to bob 1: theta1 = 35 deg from vertical down
% bob1 at (L1 sin t1, -L1 cos t1), L1 = 3
\coordinate (b1) at ({3*sin(35)}, {-3*cos(35)});
% downward vertical guide from pivot
\draw[vert] (0,0) -- (0,-3.4);
\draw[link] (0,0) -- (b1);
\node[bob] (B1) at (b1) {};
\node[label, anchor=west] at (b1) {$m_1$};
% angle arc theta1
\draw[ang] (0,-1.0) arc (270:{270+35}:1.0);
\node[engred, font=\scriptsize] at ({0.75*sin(17.5)}, {-1.25*cos(17.5)}) {$\theta_1$};
% upper link length label
\node[font=\scriptsize, enggray, anchor=east] at ({1.5*sin(35)-0.15}, {-1.5*cos(35)}) {$L_1$};

% Lower link from bob1: theta2 = -20 deg from vertical down at b1, L2 = 2.6
\coordinate (b2) at ({3*sin(35) + 2.6*sin(-20)}, {-3*cos(35) - 2.6*cos(-20)});
\draw[vert] (b1) -- ($(b1)+(0,-3.0)$);
\draw[link] (b1) -- (b2);
\node[bob] (B2) at (b2) {};
\node[label, anchor=west] at (b2) {$m_2$};
% angle arc theta2 (negative, to the left of vertical)
\draw[ang] ($(b1)+(0,-0.9)$) arc (270:{270-20}:0.9);
\node[engred, font=\scriptsize] at ($(b1)+({-0.7*sin(10)},{-1.15})$) {$\theta_2$};
\node[font=\scriptsize, enggray, anchor=west] at ($(b1)!0.5!(b2)+(0.12,0)$) {$L_2$};

\end{tikzpicture}
\caption[Double-pendulum geometry and generalised coordinates]{The
planar double pendulum and its two generalised coordinates. Each swing
angle $\theta_1$ and $\theta_2$ is measured from the downward vertical
(grey dashed) at its own pivot: $\theta_1$ at the fixed pivot, $\theta_2$
at the first bob. The four Cartesian coordinates of the two bobs are
reduced to two by the inextensible-link constraints, which the angle
coordinates satisfy identically, so the Lagrangian carries no
constraint forces. The two coordinates make the system minimal: any
fewer cannot locate it, any more would be redundant. The chapter's
project derives the equations of motion in these coordinates two ways
and reconciles them.}
\label{fig:vol03ch05:double-pendulum}
\end{figure}


\textbf{Cartesian positions of the bobs.}
\begin{align*}
(x_1, y_1) &= (L_1 \sin\theta_1, -L_1 \cos\theta_1), \\
(x_2, y_2) &= (L_1 \sin\theta_1 + L_2 \sin\theta_2, -L_1 \cos\theta_1
- L_2 \cos\theta_2).
\end{align*}

\textbf{Kinetic energy.} Differentiating the positions and computing
$\tfrac{1}{2}m_1(\dot{x}_1^2 + \dot{y}_1^2) + \tfrac{1}{2}m_2(\dot{x}_2^2
+ \dot{y}_2^2)$ gives, after simplification,
\[
T = \tfrac{1}{2}(m_1 + m_2) L_1^2 \dot{\theta}_1^2 + \tfrac{1}{2}
m_2 L_2^2 \dot{\theta}_2^2 + m_2 L_1 L_2 \dot{\theta}_1 \dot{\theta}_2
\cos(\theta_1 - \theta_2).
\]

\textbf{Potential energy.}
\[
U = -(m_1 + m_2) g L_1 \cos\theta_1 - m_2 g L_2 \cos\theta_2 + \text{
const}.
\]

\textbf{Euler-Lagrange equations.} With $L = T - U$, the two
equations are
\begin{align*}
& (m_1 + m_2) L_1^2 \ddot{\theta}_1 + m_2 L_1 L_2 \ddot{\theta}_2 \cos
(\theta_1 - \theta_2) + m_2 L_1 L_2 \dot{\theta}_2^2 \sin(\theta_1 -
\theta_2) \\
& \quad + (m_1 + m_2) g L_1 \sin\theta_1 = 0, \\
& m_2 L_2^2 \ddot{\theta}_2 + m_2 L_1 L_2 \ddot{\theta}_1 \cos(
\theta_1 - \theta_2) - m_2 L_1 L_2 \dot{\theta}_1^2 \sin(\theta_1 -
\theta_2) \\
& \quad + m_2 g L_2 \sin\theta_2 = 0.
\end{align*}

The pair is coupled, nonlinear, and analytically intractable except
in special cases (small amplitudes, in-phase motion, the limit
$m_1 \to \infty$ or $m_1 \to 0$). The general motion is chaotic in
the technical sense: the system is deterministic but exhibits
exponential sensitivity to initial conditions. The double pendulum
is the simplest mechanical demonstration of deterministic chaos and
is the project of this chapter.

The Newtonian derivation of these same equations requires drawing two
free-body diagrams, identifying the two unknown tensions (one in each
string), resolving each tension into Cartesian or polar components,
and eliminating the tensions from the coupled equations of motion. The
algebra is several pages and prone to sign errors. The Lagrangian
derivation is the two-page computation above, with no constraint
forces to eliminate.

\textbf{Small-amplitude normal modes.} For equal masses $m_1 = m_2 = m$
and equal lengths $L_1 = L_2 = L$, the small-amplitude limit
($\theta_i \ll 1$, $\cos(\theta_1-\theta_2)\to 1$, $\sin\theta_i\to
\theta_i$, products of velocities dropped) linearises the pair to
\[
2\ddot{\theta}_1 + \ddot{\theta}_2 + 2\,\tfrac{g}{L}\,\theta_1 = 0,
\qquad
\ddot{\theta}_1 + \ddot{\theta}_2 + \tfrac{g}{L}\,\theta_2 = 0.
\]
Seeking solutions $\theta_i = a_i \cos\omega t$ turns this into the
generalised eigenvalue problem $(K - \omega^2 M)\,a = 0$, with
\[
M = \begin{pmatrix} 2 & 1 \\ 1 & 1 \end{pmatrix},
\qquad
K = \frac{g}{L}\begin{pmatrix} 2 & 0 \\ 0 & 1 \end{pmatrix}.
\]
The two roots are $\omega^2 = (2 \mp \sqrt{2})\,g/L$, that is
\[
\omega_1^2 = (2 - \sqrt2)\,\tfrac{g}{L} \approx 0.586\,\tfrac{g}{L},
\qquad
\omega_2^2 = (2 + \sqrt2)\,\tfrac{g}{L} \approx 3.414\,\tfrac{g}{L}.
\]
The low-frequency mode has the two bobs swinging in phase
($a_2/a_1 = +\sqrt2/2 \approx 0.707$); the high-frequency mode has them
out of phase ($a_2/a_1 = -\sqrt2 \approx -1.414$). The file
\path{data/double_pendulum_modes.csv} carries the two modes' frequency
factors and amplitude ratios. The double pendulum is thus orderly at
small amplitude, a pair of clean normal modes, and chaotic at large
amplitude, where the linearisation that produced the modes no longer
holds.

\textbf{Numerical exploration and sensitivity.} The full nonlinear pair
has no closed-form solution and is integrated numerically. The
script \path{code/double_pendulum.py} solves the equations of motion
with a fixed-step fourth-order Runge-Kutta integrator, inverting the
two-by-two mass matrix at each step to extract the angular
accelerations, and logs the total mechanical energy as the
accuracy check: a correct conservative integration holds the energy
constant to within the integrator's truncation error, and the run
recorded in \path{data/double_pendulum_trajectory.csv} drifts by under
$2\times 10^{-7}$ in relative terms over thirty seconds. Running the
same initial condition twice with the upper angle perturbed by
$10^{-6}\,\si{\radian}$ produces trajectories whose separation grows
exponentially until it saturates at order one radian, the signature of
deterministic chaos plotted in
figure~\ref{fig:vol03ch05:double-pendulum-divergence}. The equations
are exact and the system is deterministic; it is also unpredictable in
practice past a few seconds, which is the distinction between
determinism and predictability that the chaos literature makes
central.

% Figure: sensitivity of the double pendulum. Two trajectories with
% initial conditions differing by 1e-3 rad diverge after a few seconds;
% the angle difference grows roughly exponentially then saturates.
\begin{figure}[htbp]
\centering
\begin{tikzpicture}[
  axis/.style={draw=engblack, -{Stealth}, thick},
  growth/.style={draw=engred, very thick},
  envel/.style={draw=engblue, very thick, dashed},
  guide/.style={draw=enggray, dashed, thin},
  label/.style={font=\small},
]
% x = time 0..12 s (plotted 0..10), y = log10 of angle difference, -6..1 (plotted 0..5)
\draw[axis] (0,0) -- (10.6,0) node[anchor=west, label] {$t$ [s]};
\draw[axis] (0,0) -- (0,5.6)  node[anchor=south, label] {$\log_{10}|\Delta\theta_2|$};

% x ticks (1.2 s per unit)
\foreach \x/\m in {0/0, 2.5/3, 5/6, 7.5/9, 10/12} {
  \draw (\x,-0.08) -- (\x,0.08);
  \node[font=\scriptsize, below] at (\x,-0.12) {\m};
}
% y ticks: log10 from -6 (y=0) to +1 (y=5); step 1.4 per decade
\foreach \y/\h in {0/$-6$, 1.43/$-4$, 2.86/$-2$, 4.29/$0$} {
  \draw (-0.08,\y) -- (0.08,\y);
  \node[font=\scriptsize, left] at (-0.12,\y) {\h};
}

% Linear-in-log growth = exponential divergence, slope = Lyapunov exponent
% start at log10 = -6 (separation 1e-6 in theta1 maps to ~1e-6 in theta2 after coupling)
% reach saturation ~ log10 = 0.3 (order 1 rad) by t ~ 9 s, then flatten
% slope: from (0,0) to (7.5, 4.3): straight line
\draw[growth] (0,0) -- (7.0, 4.3);
% saturation plateau
\draw[growth] (7.0,4.3) .. controls (8.0,4.6) and (9.0,4.55) .. (10,4.5);
\node[engred, label, anchor=south west] at (2.6, 1.6) {$|\Delta\theta_2|\sim e^{\lambda t}$};

% slope guide annotation
\draw[guide] (1,0.61) -- (3,0.61) -- (3, {0.61 + 2*0.61});
\node[font=\scriptsize, enggray, anchor=west] at (3.1, 1.0) {slope $=\lambda/\ln 10$};

% saturation note
\node[engblue, font=\scriptsize, anchor=south east] at (10, 4.8) {saturation $\sim O(1)$ rad};
\draw[guide] (0, 4.5) -- (10, 4.5);

\end{tikzpicture}
\caption[Exponential divergence of nearby double-pendulum
trajectories]{The signature of deterministic chaos in the double
pendulum. Two simulations start from initial conditions differing by
$10^{-6}\,\si{\radian}$ in the upper angle. The base-ten logarithm of
the resulting difference $|\Delta\theta_2(t)|$ grows linearly with time
(red), which means the difference itself grows exponentially,
$|\Delta\theta_2|\sim e^{\lambda t}$; the slope divided by $\ln 10$ is
the largest Lyapunov exponent $\lambda$. Growth continues until the two
trajectories are fully uncorrelated, where the separation saturates at
order one radian (blue dashed). The energy method delivers the exact
equations of motion, and those exact equations are unpredictable in
practice: determinism and predictability are not the same property.}
\label{fig:vol03ch05:double-pendulum-divergence}
\end{figure}


\subsection{Bead on a rotating hoop}

A bead of mass $m$ slides without friction on a vertical circular
hoop of radius $R$ that rotates about its vertical diameter with
constant angular velocity $\Omega$. Find the bead's stable
equilibrium positions on the hoop as a function of $\Omega$.

\textbf{Generalised coordinate.} The angle $\theta$ measuring the
bead's position on the hoop from the bottom; $\theta = 0$ is the
lowest point, $\theta = \pi$ the highest.

\textbf{Kinetic energy.} The bead's velocity has two components: a
tangential velocity $R\dot{\theta}$ along the hoop, and a tangential
velocity $R\sin\theta \cdot \Omega$ due to the hoop's rotation. The
two components are perpendicular, so
\[
T = \tfrac{1}{2} m (R\dot{\theta})^2 + \tfrac{1}{2} m (R\sin\theta
\cdot \Omega)^2 = \tfrac{1}{2} m R^2 (\dot{\theta}^2 + \Omega^2
\sin^2\theta).
\]

\textbf{Potential energy.} The bead's height above the bottom is
$R(1 - \cos\theta)$, so $U = m g R (1 - \cos\theta)$.

\textbf{Lagrangian and equation of motion.} The Lagrangian $L = T - U$
gives
\[
m R^2 \ddot{\theta} - m R^2 \Omega^2 \sin\theta \cos\theta + m g R
\sin\theta = 0.
\]
At equilibrium $\ddot{\theta} = 0$, so
\[
\sin\theta \,(g - R \Omega^2 \cos\theta) = 0.
\]
The solutions are $\theta = 0$ (bottom), $\theta = \pi$ (top), and,
when $R\Omega^2 \cos\theta = g$ has a solution in $(0, \pi/2)$, two
additional equilibria $\theta = \pm \theta_c$ with $\cos\theta_c =
g/(R\Omega^2)$.

The additional equilibria exist only when $R\Omega^2 > g$, i.e.,
when the centripetal acceleration at the equator of the hoop exceeds
gravity. Below the threshold, the bead's stable position is at the
bottom; above the threshold, the bottom becomes unstable and the
bead settles at $\pm \theta_c$. The transition at $R\Omega^2 = g$ is
a \engterm{bifurcation}, the simplest mechanical example of a
symmetry-breaking transition; it is the prototype for many phase
transitions in continuous physical systems. The branch diagram in
figure~\ref{fig:vol03ch05:hoop-bifurcation} draws the pitchfork: the
bottom equilibrium loses stability at the threshold and two symmetric
off-axis equilibria appear and take it over.

The stability of each equilibrium is read directly off an effective
potential. Collecting the position-dependent terms of the Lagrangian
into $U_\text{eff}(\theta) = m g R(1 - \cos\theta) - \tfrac12 m R^2
\Omega^2 \sin^2\theta$, the equilibria are its stationary points and the
stable ones are its minima, identified by the sign of $U_\text{eff}''$.
Below the threshold the bottom is the sole minimum; above it the bottom
becomes a maximum and the two off-axis points become the minima. The
rotation term has reshaped the potential, lowering the energy of
off-axis positions until they out-compete the bottom. The script
\path{code/hoop_bifurcation.py} tabulates the equilibrium angles and
their stability across the threshold into
\path{data/hoop_equilibria.csv}, and the symbolic derivation of the
equation of motion is in \path{code/lagrangian_symbolic.py}, which
generates the Euler-Lagrange equation for this system and for the
simple pendulum from the Lagrangian alone, with no hand algebra.

% Figure: bifurcation diagram for the bead on a rotating hoop. Below the
% critical speed the bottom is the only stable equilibrium; above it the
% bottom destabilises and two symmetric off-axis equilibria appear.
\begin{figure}[htbp]
\centering
\begin{tikzpicture}[
  axis/.style={draw=engblack, -{Stealth}, thick},
  stable/.style={draw=engblue, very thick},
  unstable/.style={draw=engred, very thick, dashed},
  guide/.style={draw=enggray, dashed, thin},
  label/.style={font=\small},
]
% x = rotation parameter R Omega^2 / g from 0 to 3 (plotted 0..9)
% y = equilibrium angle theta_eq from -pi/2 to pi/2 (plotted -3..3)
\draw[axis] (0,0) -- (9.6,0) node[anchor=west, label] {$R\Omega^2/g$};
\draw[axis] (0,-3.4) -- (0,3.4) node[anchor=north west, label] {$\theta_\text{eq}$};

% critical value at R Omega^2 / g = 1 -> x = 3
\draw[guide] (3,-3.2) -- (3,3.2);
\node[font=\scriptsize, enggray, anchor=south] at (3,3.0) {threshold $R\Omega^2=g$};

% x ticks
\foreach \x/\m in {3/$1$, 6/$2$, 9/$3$} {
  \draw (\x,-0.08) -- (\x,0.08);
  \node[font=\scriptsize, below] at (\x,-0.12) {\m};
}
% y ticks
\node[font=\scriptsize, left] at (-0.12, 0) {$0$};
\node[font=\scriptsize, left] at (-0.12, 3) {$\pi/2$};
\node[font=\scriptsize, left] at (-0.12, -3) {$-\pi/2$};

% Branch theta = 0: stable for x < 3, unstable for x > 3
\draw[stable] (0,0) -- (3,0);
\draw[unstable] (3,0) -- (9,0);
\node[engblue, font=\scriptsize, anchor=south] at (1.5, 0.05) {bottom stable};
\node[engred, font=\scriptsize, anchor=north] at (6, -0.05) {bottom now unstable};

% Off-axis branches: cos(theta_c) = g/(R Omega^2) = 1/p, p = x/3
% theta_c = acos(3/x) for x > 3, in radians; plotted y = theta_c * (3 / (pi/2)) = theta_c * 1.9099
% sample at x = 3..9
\draw[stable] plot[domain=3.05:9, samples=60]
  ({\x}, {1.9099 * acos(3/\x) * 3.14159/180});
\draw[stable] plot[domain=3.05:9, samples=60]
  ({\x}, {-1.9099 * acos(3/\x) * 3.14159/180});
\node[engblue, font=\scriptsize, anchor=south west] at (6.5, 2.0) {$+\theta_c$ stable};
\node[engblue, font=\scriptsize, anchor=north west] at (6.5, -2.0) {$-\theta_c$ stable};

\end{tikzpicture}
\caption[Pitchfork bifurcation of the bead on a rotating hoop]{The
equilibrium angle of a bead on a frictionless vertical hoop rotating at
angular speed $\Omega$, as a function of the dimensionless group
$R\Omega^2/g$. Below the threshold $R\Omega^2=g$ the bottom of the hoop
($\theta=0$) is the only stable equilibrium (solid blue). At the
threshold the bottom loses stability (red dashed) and two symmetric
off-axis equilibria $\pm\theta_c$ with $\cos\theta_c=g/(R\Omega^2)$
branch off and are stable. The shape is the pitchfork bifurcation, the
prototype of symmetry-breaking transitions; the energy method delivers
it from a single line, the stationarity of an effective potential.}
\label{fig:vol03ch05:hoop-bifurcation}
\end{figure}


\section{Strain energy and the deflection of structures}\pathtag{standard}

The dynamics of sections 5.3 through 5.6 use kinetic and potential
energy of moving masses. A second, equally large branch of energy
methods works on structures at rest: a loaded elastic member stores
\engterm{strain energy} as it deforms, and the deflection of the
structure under load can be recovered from that scalar by
differentiation, with no need to integrate the deflected shape or solve
the governing differential equation. The two branches share the same
discipline. Write one scalar energy, differentiate it in a standard
pattern, and read the quantity of interest off the result.

\engterm{Strain energy} is the elastic potential energy stored in a
deformed member, equal to the work done by the loads in deforming it.
For an axially loaded bar of length $L$, cross-sectional area $A$, and
modulus $E$ carrying axial force $N$, the stored energy is
\[
U = \frac{N^2 L}{2 E A},
\]
the area under the linear force-extension line, $\tfrac12 N \delta$ with
$\delta = NL/(EA)$. For a beam in bending with internal bending moment
$M(x)$, second moment of area $I$, and modulus $E$, the strain energy is
the integral of the bending contribution along the member,
\[
U = \int_0^L \frac{M(x)^2}{2 E I}\,dx,
\]
with parallel expressions for torsion ($T^2 L / (2GJ)$ per uniform shaft)
and for transverse shear (usually small in slender members and dropped
here). Each is a quadratic, positive form in the internal load, the
structural analogue of the $\tfrac12 k x^2$ spring well of
section 5.1.

\engterm{Castigliano's second theorem} turns strain energy into a
deflection. For a linearly elastic structure, the displacement of the
point of application of a load, measured in the direction of that load,
equals the partial derivative of the total strain energy with respect
to the load:
\[
\delta_i = \frac{\partial U}{\partial P_i},
\qquad
\theta_i = \frac{\partial U}{\partial M_i}
\]
for a deflection conjugate to a force $P_i$ or a rotation conjugate to a
couple $M_i$. The theorem is the static, structural sibling of the
cyclic-coordinate result of section 5.3: a derivative of a scalar energy
yields a kinematic quantity. Carlo Alberto Castigliano stated it in his
1879 \emph{Th\'eorie de l'\'equilibre des syst\`emes \'elastiques}; the
working derivation and a catalogue of applications are in any
mechanical-design reference \cite[ch.~4]{text:shigley-mechanical-design}.

\subsection{End-to-end case study: tip deflection of a bent cantilever
frame}

A complete energy-method analysis, carried from geometry to a checked
number, fixes the procedure better than a statement of the theorem. The
structure is the L-shaped cantilever frame of
figure~\ref{fig:vol03ch05:castigliano-frame}: a horizontal member of
length $a$ joined at a rigid corner to a vertical member of length $b$
built into a wall, with a vertical load $P$ at the free tip. Both members
are steel tubes of the same uniform bending stiffness $EI$. We want the
vertical deflection of the loaded tip.

% Figure: an L-shaped (bent) cantilever frame loaded by a vertical force P
% at the free tip. The figure fixes the geometry and the two segments over
% which the strain-energy integrals run for the Castigliano case study.
\begin{figure}[htbp]
\centering
\begin{tikzpicture}[
  member/.style={draw=engblack, very thick},
  ghost/.style={draw=enggray, dashed, thick},
  force/.style={draw=engred, -{Stealth}, very thick},
  dimline/.style={draw=engblue, {Stealth}-{Stealth}, thin},
  label/.style={font=\small},
]
% Fixed wall at top-left, vertical member BC down, horizontal member CA to the right.
% B = built-in support (top), C = corner, A = free tip with load P.
\coordinate (B) at (0,4);
\coordinate (C) at (0,0);
\coordinate (A) at (4,0);

% hatched wall at the built-in end B
\draw[member] (-0.5,4) -- (0.5,4);
\foreach \x in {-0.4,-0.2,0,0.2,0.4} {
  \draw[draw=enggray, thin] (\x,4) -- (\x-0.18,4.28);
}
\node[font=\scriptsize, anchor=south] at (0,4.34) {built-in};

% vertical member B==C (length b), horizontal member C==A (length a)
\draw[member] (B) -- (C);
\draw[member] (C) -- (A);

% corner and tip markers
\fill[engblack] (C) circle (2pt);
\node[font=\scriptsize, anchor=north east] at (C) {$C$};
\fill[engblack] (A) circle (2pt);
\node[font=\scriptsize, anchor=north] at (4,-0.15) {$A$};
\node[font=\scriptsize, anchor=south east] at (0,3.9) {$B$};

% applied load P at the free tip A, pointing down
\draw[force] (4,1.3) -- (4,0.12);
\node[engred, label, anchor=west] at (4.05,0.8) {$P$};

% deflected shape (exaggerated) showing the tip dropping
\draw[ghost] (0,4) .. controls (0.5,2) and (1.2,0.6) .. (4,-0.85);

% dimension a (horizontal, along C==A)
\draw[dimline] (0,-0.7) -- (4,-0.7);
\node[engblue, font=\scriptsize, anchor=north] at (2,-0.72) {$a$};
% dimension b (vertical, along B==C)
\draw[dimline] (-0.85,0) -- (-0.85,4);
\node[engblue, font=\scriptsize, anchor=east] at (-0.88,2) {$b$};

% local coordinate hints for the strain-energy integrals
\node[engamber, font=\scriptsize, anchor=south] at (2,0.12) {$x:\,0\to a$};
\node[engamber, font=\scriptsize, anchor=west] at (0.1,2) {$y:\,0\to b$};

\end{tikzpicture}
\caption[Bent cantilever frame for the Castigliano case study]{The
L-shaped cantilever frame of the worked case study: a horizontal member
$CA$ of length $a$ joined at a rigid corner $C$ to a vertical member $BC$
of length $b$ built in at $B$. A vertical load $P$ at the free tip $A$
bends both members. The bending moment in the horizontal member is $M(x)
= P x$ for $x$ measured from the tip; in the vertical member it is the
constant $M = P a$, since the whole load $P$ acts at lever arm $a$ from
every point on $BC$. Castigliano's second theorem reads the tip
deflection off these two moment expressions by differentiating the total
strain energy with respect to $P$. The grey curve is the deflected shape,
drawn to exaggerated scale.}
\label{fig:vol03ch05:castigliano-frame}
\end{figure}


\textbf{Step 1: internal bending moments.} Walk the structure from the
free end inward, since the load and geometry are simplest there. On the
horizontal member, measure $x$ from the tip toward the corner. The only
load outboard of the cut is $P$ at lever arm $x$, so the bending moment
is
\[
M_1(x) = P x,\qquad 0 \le x \le a.
\]
On the vertical member, measure $y$ from the corner up toward the wall.
Every point on this member sees the full load $P$ acting at the constant
horizontal lever arm $a$ (the horizontal distance from the member to the
tip does not change as $y$ increases), so
\[
M_2(y) = P a,\qquad 0 \le y \le b,
\]
a constant. The axial force in each member (the horizontal member carries
no axial load from a vertical tip force; the vertical member carries the
direct compression $P$) contributes negligibly to the deflection of a
slender frame and we drop it, returning to it as a check at the end.

\textbf{Step 2: total strain energy.} Sum the bending contributions of
the two members,
\[
U = \int_0^a \frac{M_1(x)^2}{2EI}\,dx
  + \int_0^b \frac{M_2(y)^2}{2EI}\,dy
  = \int_0^a \frac{(Px)^2}{2EI}\,dx
  + \int_0^b \frac{(Pa)^2}{2EI}\,dy.
\]
The first integral is $\dfrac{P^2}{2EI}\cdot\dfrac{a^3}{3}$; the second,
with a constant integrand, is $\dfrac{P^2 a^2}{2EI}\cdot b$. So
\[
U = \frac{P^2 a^3}{6EI} + \frac{P^2 a^2 b}{2EI}
  = \frac{P^2 a^2}{2EI}\left(\frac{a}{3} + b\right).
\]

\textbf{Step 3: differentiate for the deflection.} The tip deflection in
the direction of $P$ is Castigliano's derivative,
\[
\delta_\text{tip} = \frac{\partial U}{\partial P}
  = \frac{P a^2}{EI}\left(\frac{a}{3} + b\right)
  = \frac{P a^3}{3EI} + \frac{P a^2 b}{EI}.
\]
The two terms have a clean reading. The first is the deflection of a
plain cantilever of length $a$ under a tip load, $Pa^3/(3EI)$, the
standard result. The second is the extra tip drop caused by the vertical
member bending under the constant moment $Pa$; it scales with $b$ because
a longer vertical member rotates the corner through a larger angle, and
that rotation throws the tip down by the moment arm $a$.

\textbf{Step 4: a number.} Take $a = 1.2\,\si{\meter}$,
$b = 0.8\,\si{\meter}$, a load $P = 2.0\,\si{\kilo\newton}$, and steel
tubing with $E = 200\,\si{\giga\pascal}$ and
$I = 4.0\times 10^{-6}\,\si{\meter\tothe{4}}$ (a $60\,\si{\milli\meter}$
square tube of $\sim 4\,\si{\milli\meter}$ wall, say), so
$EI = 8.0\times 10^{5}\,\si{\newton\meter\squared}$. Then
\[
\delta_\text{tip}
= \frac{2000 \cdot 1.2^2}{8.0\times 10^{5}}
  \left(\frac{1.2}{3} + 0.8\right)
= \frac{2000 \cdot 1.44}{8.0\times 10^{5}} \cdot 1.2
\approx 4.3\times 10^{-3}\,\si{\meter},
\]
a tip deflection of about $4.3\,\si{\milli\meter}$. The cantilever term
alone supplies $Pa^3/(3EI) = 2000\cdot 1.728/(2.4\times 10^{6}) \approx
1.44\,\si{\milli\meter}$; the corner-rotation term adds the remaining
$\sim 2.9\,\si{\milli\meter}$. The vertical member, though it carries
only a constant moment, contributes two-thirds of the tip movement, a
result a cantilever-only estimate would miss by a factor of three.

\textbf{Step 5: cross-check by a direct method.} The same deflection
follows from the moment-area or unit-load route without Castigliano, and
agreement between two independent methods is the discipline's standard
verification. Apply a virtual unit vertical load at the tip; its bending
moments are $m_1(x) = x$ on the horizontal member and $m_2(y) = a$ on the
vertical member, the real moments with $P$ set to $1$. The unit-load
deflection integral is
\[
\delta_\text{tip}
= \int_0^a \frac{M_1 m_1}{EI}\,dx + \int_0^b \frac{M_2 m_2}{EI}\,dy
= \int_0^a \frac{(Px)(x)}{EI}\,dx + \int_0^b \frac{(Pa)(a)}{EI}\,dy
= \frac{P a^3}{3EI} + \frac{P a^2 b}{EI},
\]
identical to the Castigliano result term for term. Two methods, one
number; the analysis is closed. The dropped axial term in the vertical
member would add $PL_2/(EA)$ at the order of $P b/(EA) \approx 2000 \cdot
0.8/(200\times 10^9 \cdot 9\times 10^{-4}) \approx 9\times 10^{-9}\,
\si{\meter}$, six orders of magnitude below the bending deflection, which
justifies dropping it for this slender frame.

\subsection{Worked example: deflection of a truss joint by the unit-load
method}

The unit-load method that cross-checked the frame is the standard tool
for the deflection of a pin-jointed truss, where the members carry only
axial force and the bending integrals collapse to a sum over bars. The
deflection of a chosen joint in a chosen direction is
\[
\delta = \sum_i \frac{n_i N_i L_i}{E_i A_i},
\]
with $N_i$ the real axial force in bar $i$ under the actual load, $n_i$
the axial force in the same bar under a virtual unit load applied at the
joint in the wanted direction, and $L_i/(E_i A_i)$ the bar's axial
flexibility. Consider the two-bar truss of
figure~\ref{fig:vol03ch05:truss-unit-load}: a horizontal bar $BC$ of
length $a = 3.0\,\si{\meter}$ and a diagonal bar $AC$ running from a
higher wall pin $A$ down to the loaded joint $C$, with vertical rise
$h = 2.25\,\si{\meter}$, so the diagonal length is $L_{AC} =
\sqrt{a^2 + h^2} = 3.75\,\si{\meter}$. A vertical load $P =
15\,\si{\kilo\newton}$ hangs at $C$. Both bars are steel with
$EA = 4.0\times 10^{7}\,\si{\newton}$.

% Figure: a simple three-bar pin-jointed truss. A vertical load P at the
% loaded joint; the unit-load method computes the vertical deflection of
% that joint from the real bar forces N_i and the unit-load bar forces n_i.
\begin{figure}[htbp]
\centering
\begin{tikzpicture}[
  member/.style={draw=engblack, very thick},
  force/.style={draw=engred, -{Stealth}, very thick},
  unit/.style={draw=englime, -{Stealth}, very thick},
  pin/.style={draw=engblack, fill=white, circle, inner sep=1.8pt},
  label/.style={font=\small},
]
% Two wall supports at left (A top, B bottom) and a free joint C to the right.
\coordinate (A) at (0,3);
\coordinate (B) at (0,0);
\coordinate (C) at (4,0);

% hatched wall
\draw[member] (-0.4,-0.4) -- (-0.4,3.4);
\foreach \y in {-0.2,0.2,0.6,1.0,1.4,1.8,2.2,2.6,3.0} {
  \draw[draw=enggray, thin] (-0.4,\y) -- (-0.7,\y+0.2);
}

% members: A==C (diagonal), B==C (horizontal)
\draw[member] (A) -- (C);
\draw[member] (B) -- (C);

% pins
\node[pin] at (A) {};
\node[pin] at (B) {};
\node[pin] at (C) {};
\node[font=\scriptsize, anchor=east] at (-0.05,3) {$A$};
\node[font=\scriptsize, anchor=east] at (-0.05,0) {$B$};
\node[font=\scriptsize, anchor=south west] at (4.05,0.05) {$C$};

% member labels with lengths
\node[font=\scriptsize, anchor=south, engblack] at (2.0,1.6) {bar 1 ($AC$)};
\node[font=\scriptsize, anchor=north, engblack] at (2.0,-0.05) {bar 2 ($BC$)};

% geometry dimensions
\node[engblue, font=\scriptsize, anchor=east] at (-0.05,1.5) {$h$};
\draw[draw=engblue, {Stealth}-{Stealth}, thin] (0,-0.55) -- (4,-0.55);
\node[engblue, font=\scriptsize, anchor=north] at (2,-0.57) {$a$};

% real load P (red) downward at C
\draw[force] (4,1.5) -- (4,0.18);
\node[engred, label, anchor=west] at (4.05,0.9) {$P$};

% virtual unit load (green) downward at C, drawn offset to the side
\draw[unit] (5.0,1.5) -- (5.0,0.6);
\node[englime, font=\scriptsize, anchor=west] at (5.05,1.05) {unit load $1$};
\node[englime, font=\scriptsize, anchor=west] at (5.05,0.45) {(virtual)};

\end{tikzpicture}
\caption[Pin-jointed truss for the unit-load method]{A two-bar
pin-jointed truss carrying a vertical load $P$ at the free joint $C$:
a diagonal bar $AC$ and a horizontal bar $BC$, both pinned to a rigid
wall at $A$ and $B$. The unit-load (or virtual-work) method finds the
vertical deflection of joint $C$ in two passes. First, statics gives the
real axial force $N_i$ in each bar under the load $P$. Second, a virtual
unit load (green) is applied at $C$ in the direction of the wanted
deflection, and statics gives the virtual axial force $n_i$ in each bar.
The deflection is then the sum $\delta_C = \sum_i n_i N_i L_i /(E_i A_i)$
over the bars, the strain-energy pairing of real and virtual force
systems. The method computes a single deflection component without
solving for the whole deformed geometry.}
\label{fig:vol03ch05:truss-unit-load}
\end{figure}


\textbf{Real bar forces.} Resolve at joint $C$. The diagonal makes an
angle $\phi$ with the horizontal where $\sin\phi = h/L_{AC} = 2.25/3.75 =
0.6$ and $\cos\phi = a/L_{AC} = 3.0/3.75 = 0.8$. Vertical equilibrium at
$C$ requires the diagonal to carry the load: $N_{AC}\sin\phi = P$, so
$N_{AC} = P/0.6 = 25\,\si{\kilo\newton}$ (tension). Horizontal
equilibrium requires the horizontal bar to balance the diagonal's
horizontal pull: $N_{BC} = -N_{AC}\cos\phi = -25\cdot 0.8 =
-20\,\si{\kilo\newton}$ (compression).

\textbf{Virtual unit-load forces.} A virtual unit downward load at $C$
produces, by the identical statics with $P \to 1$, $n_{AC} = 1/0.6 =
1.667$ and $n_{BC} = -1.333$.

\textbf{Assemble the deflection.}
\[
\delta_C = \frac{n_{AC} N_{AC} L_{AC}}{EA}
         + \frac{n_{BC} N_{BC} L_{BC}}{EA}.
\]
The two products are $n_{AC}N_{AC} = 1.667\cdot 25000 = 41{,}670\,
\si{\newton}$ over length $3.75\,\si{\meter}$, and $n_{BC}N_{BC} =
(-1.333)(-20000) = 26{,}670\,\si{\newton}$ over length
$3.0\,\si{\meter}$. Both products are positive: the compression bar
contributes a positive deflection because both its real and virtual
forces are compressive and their product is positive. So
\[
\delta_C = \frac{41{,}670 \cdot 3.75 + 26{,}670 \cdot 3.0}
                {4.0\times 10^{7}}
= \frac{156{,}300 + 80{,}000}{4.0\times 10^{7}}
\approx 5.9\times 10^{-3}\,\si{\meter},
\]
a downward joint deflection of about $5.9\,\si{\milli\meter}$. The method
delivered one deflection component from two statics passes and a sum over
two bars, with no deformed-geometry sketch and no compatibility equation.

\subsection{End-to-end case study: energy absorption in a drop impact}

Energy methods carry their largest practical weight in impact, where the
peak force a structure must survive is many times the static weight of
the striking mass and where a direct force analysis would need the full
time history of the collision. The energy route sidesteps the time
history: equate the work the falling mass does to the strain energy the
structure stores at peak deflection, and the peak force and deflection
follow from one scalar balance.

\textbf{The setup.} A mass $m = 50\,\si{\kilogram}$ is dropped from a
height $h = 0.30\,\si{\meter}$ onto the free tip of a steel cantilever of
length $L = 1.0\,\si{\meter}$ and bending stiffness $EI = 3.0\times
10^{4}\,\si{\newton\meter\squared}$ (a light beam, chosen so the impact
is severe). Find the peak deflection, the peak force on the beam, and the
dynamic load factor by which the impact exceeds a static placement of the
same mass.

\textbf{Static reference.} Placed gently at the tip, the weight $W = mg =
50 \cdot 9.81 = 491\,\si{\newton}$ produces the static tip deflection of
a cantilever,
\[
\delta_\text{st} = \frac{W L^3}{3 EI}
= \frac{491 \cdot 1.0^3}{3 \cdot 3.0\times 10^{4}}
\approx 5.45\times 10^{-3}\,\si{\meter},
\]
about $5.5\,\si{\milli\meter}$. The tip behaves as a linear spring of
stiffness $k = W/\delta_\text{st} = 3EI/L^3 = 9.0\times 10^{4}\,
\si{\newton\per\meter}$.

\textbf{Energy balance.} The mass falls the drop height $h$ plus the
extra distance $\delta_\text{max}$ it pushes the tip down before
stopping, doing total work $W(h + \delta_\text{max})$. At peak deflection
the mass is momentarily at rest, so all of that work is stored as strain
energy $\tfrac12 k\,\delta_\text{max}^2$ in the beam:
\[
W(h + \delta_\text{max}) = \tfrac12 k\,\delta_\text{max}^2.
\]
Writing the static deflection $\delta_\text{st} = W/k$ and rearranging
gives a quadratic in $\delta_\text{max}$ whose physical root is
\[
\delta_\text{max} = \delta_\text{st}\left(1 + \sqrt{1 +
\frac{2h}{\delta_\text{st}}}\right) = n\,\delta_\text{st},
\]
with the dimensionless \engterm{dynamic load factor}
\[
n = 1 + \sqrt{1 + \frac{2h}{\delta_\text{st}}}.
\]
The factor depends only on the ratio of drop height to static deflection,
plotted in figure~\ref{fig:vol03ch05:impact-factor}.

% Figure: the dynamic load (impact) factor n = 1 + sqrt(1 + 2h/delta_st)
% plotted against the drop-height ratio h/delta_st. The factor measures how
% much a suddenly applied or dropped load amplifies the static response.
\begin{figure}[htbp]
\centering
\begin{tikzpicture}[
  axis/.style={draw=engblack, -{Stealth}, thick},
  curve/.style={draw=engblue, very thick},
  guide/.style={draw=enggray, dashed, thin},
  critmark/.style={draw=engred, dashed, thin},
  label/.style={font=\small},
]
% x = h/delta_st from 0 to 50 (plotted 0..10, scale 1 unit = 5 of ratio)
% y = impact factor n from 0 to 12 (plotted 0..6, scale 1 unit = 2 of n)
\draw[axis] (0,0) -- (10.6,0) node[anchor=west, label] {$h/\delta_\mathrm{st}$};
\draw[axis] (0,0) -- (0,6.4) node[anchor=south, label] {impact factor $n$};

% x ticks: ratio 0,10,20,30,40,50  -> plotted at 0,2,4,6,8,10
\foreach \xp/\xv in {0/0, 2/10, 4/20, 6/30, 8/40, 10/50} {
  \draw (\xp,-0.08) -- (\xp,0.08);
  \node[font=\scriptsize, below] at (\xp,-0.12) {$\xv$};
}
% y ticks: n 0,2,4,6,8,10,12 -> plotted at 0,1,2,3,4,5,6
\foreach \yp/\yv in {0/0, 1/2, 2/4, 3/6, 4/8, 5/10, 6/12} {
  \draw (-0.08,\yp) -- (0.08,\yp);
  \node[font=\scriptsize, left] at (-0.12,\yp) {$\yv$};
}

% n = 1 + sqrt(1 + 2 r), with r = h/delta_st.
% plotted: X in 0..10 maps to r = 5*X; Y = n/2.
% clamp the argument of sqrt: 1 + 2r is always >= 1 here, safe.
\draw[curve] plot[domain=0:10, samples=200]
  ({\x}, {(1 + sqrt(1 + 2*5*\x))/2});
\node[engblue, label, anchor=south west] at (5.2,4.1)
  {$n = 1 + \sqrt{1 + \dfrac{2h}{\delta_\mathrm{st}}}$};

% reference: sudden load (h = 0) gives n = 2
\draw[critmark] (0,1) -- (1.6,1);
\node[engred, font=\scriptsize, anchor=west] at (1.65,1) {$n=2$ (suddenly applied, $h=0$)};
\fill[engred] (0,1) circle (1.8pt);

% slowly applied (static) limit n -> 1 is the y-intercept floor
\node[enggray, font=\scriptsize, anchor=south west] at (0.1,0.05)
  {$n\to 1$ slow loading};

\end{tikzpicture}
\caption[Dynamic load (impact) factor versus drop-height ratio]{The
dynamic load factor $n = 1 + \sqrt{1 + 2h/\delta_\mathrm{st}}$ for a mass
dropped from height $h$ onto an elastic member whose static deflection
under the same weight would be $\delta_\mathrm{st}$. A load lowered slowly
gives $n\to 1$, the static case. A load released from rest at the contact
point ($h=0$) already doubles the response, $n=2$, the suddenly applied
limit. A genuine drop from even a modest height multiplies the peak force
and deflection several-fold, which is why catching a falling load on a
stiff member (small $\delta_\mathrm{st}$) is far more punishing than
catching it on a compliant one. The energy method derives the factor from
a single statement: the work done by the falling weight equals the strain
energy stored at peak deflection.}
\label{fig:vol03ch05:impact-factor}
\end{figure}


\textbf{Numbers.} The drop-height ratio is $2h/\delta_\text{st} =
2\cdot 0.30/5.45\times 10^{-3} = 110$, so
\[
n = 1 + \sqrt{1 + 110} = 1 + \sqrt{111} \approx 1 + 10.5 = 11.5.
\]
The peak deflection is $\delta_\text{max} = 11.5 \cdot 5.45\times 10^{-3}
\approx 6.3\times 10^{-2}\,\si{\meter}$, about $63\,\si{\milli\meter}$,
and the peak force is $F_\text{max} = k\,\delta_\text{max} = n\,W = 11.5
\cdot 491 \approx 5.6\,\si{\kilo\newton}$.

\textbf{Verdict.} A $50\,\si{\kilogram}$ mass that weighs half a
kilonewton when set down gently strikes the beam with more than five and
a half kilonewtons when dropped a third of a metre, an eleven-fold
amplification. A designer who sized the beam for the static weight alone
would under-rate it by a factor of eleven, and the beam would yield or
fail on the first drop. The amplification grows with the square root of
the drop height and is worst for stiff members, where $\delta_\text{st}$
is small and the ratio $2h/\delta_\text{st}$ large. The lesson the energy
method makes unavoidable is that a stiffer structure is not always a
safer one under impact: stiffness lowers the static deflection but raises
the dynamic load factor, and a compliant element deliberately added in
series (a rubber pad, a crush tube, a spring mount) lowers the peak force
by giving the impact energy a larger deflection over which to be
absorbed. This is the design principle behind crumple zones, packaging
foam, and the energy-absorbing mounts on sensitive equipment.

\begin{principle}[Impact amplifies the static response]
A load applied dynamically deflects and stresses an elastic member by a
dynamic load factor $n = 1 + \sqrt{1 + 2h/\delta_\text{st}}$ over the
static value, where $h$ is the drop height and $\delta_\text{st}$ the
static deflection under the same weight. A suddenly applied load ($h=0$)
already doubles the response. Sizing a member against the static load
alone is unsafe whenever the load can arrive with any velocity; the
energy balance between work done and strain energy stored gives the peak
without the collision's time history.
\end{principle}

\subsection{Worked example: Rayleigh estimate of a natural frequency}

The Rayleigh method of the section 5.5 mastery box estimates a vibrating
system's fundamental frequency from an energy ratio on an assumed mode
shape, the practising engineer's fastest check on a
structural-dynamics result. Estimate the
fundamental frequency of a uniform cantilever beam of length $L$, mass
per unit length $\rho A$, and bending stiffness $EI$, carrying no tip
mass.

\textbf{Trial shape.} Choose a shape that satisfies the geometric
(clamped-end) boundary conditions $\phi(0) = 0$ and $\phi'(0) = 0$. The
static-deflection shape under a uniform load is one good choice; a
simpler admissible polynomial is $\phi(x) = 1 - \cos(\pi x / (2L))$,
which gives $\phi(0) = 0$, $\phi'(0) = 0$, and a tip value $\phi(L) = 1$.

\textbf{Energy ratio.} The Rayleigh quotient for a beam is the ratio of
peak bending strain energy to a reference kinetic energy,
\[
\omega^2 = \frac{\displaystyle\int_0^L EI\,[\phi''(x)]^2\,dx}
                {\displaystyle\int_0^L \rho A\,[\phi(x)]^2\,dx}.
\]
With $\phi = 1 - \cos(\pi x/(2L))$, the second derivative is $\phi'' =
(\pi/(2L))^2\cos(\pi x/(2L))$. The numerator integral evaluates to
$EI\,(\pi/(2L))^4 \cdot (L/2)$ and the denominator to $\rho A\,L\,(3/2 -
4/\pi)$, giving
\[
\omega_1 \approx 3.66\,\sqrt{\frac{EI}{\rho A L^4}}.
\]
The exact Euler-Bernoulli coefficient is $(1.875)^2 = 3.516$, so the
crude one-term trial shape over-estimates the fundamental by about
$4\,\%$, on the right side: the Rayleigh estimate is always an
upper bound, because constraining the motion to a single assumed shape
stiffens the system. A designer who reads $4.6\,\si{\hertz}$ from a
finite-element model and computes a Rayleigh estimate of $4.8\,\si{\hertz}$
has a consistent pair; a model returning $1\,\si{\hertz}$ or
$15\,\si{\hertz}$ against this estimate carries an error in the model, not
a surprise in the structure.

\subsection{Worked example: equilibrium of a spring system by minimum
potential energy}

The static counterpart of Hamilton's principle is the principle of
\engterm{minimum potential energy}: a conservative system in stable
equilibrium sits at a minimum of its total potential energy, and the
equilibrium configuration is found by setting the energy's gradient to
zero. The method recovers force balances from a single scalar without
drawing free-body diagrams, and it extends directly to systems where a
force balance would be awkward.

A mass $m$ hangs from a spring of stiffness $k$ fixed to a support
(figure~\ref{fig:vol03ch05:min-pe-springs}). Let $x$ be the downward
displacement of the mass from the spring's natural length. The total
potential energy is the spring energy plus the gravitational energy,
\[
U(x) = \tfrac12 k x^2 - m g x,
\]
the gravitational term negative because the mass loses height as $x$
grows. The equilibrium is the stationary point,
\[
\frac{dU}{dx} = k x - m g = 0
\quad\Longrightarrow\quad
x_\text{eq} = \frac{mg}{k},
\]
the familiar static stretch, recovered without a force balance. The
second derivative $U'' = k > 0$ confirms the stationary point is a
minimum, so the equilibrium is stable; the curvature $k$ is the system's
stiffness and sets the small-oscillation frequency $\omega = \sqrt{k/m}$
about the equilibrium. The same procedure handles a mass held between two
springs of stiffnesses $k_1$ and $k_2$: the energy is $\tfrac12 k_1 x^2 +
\tfrac12 k_2 x^2 - F x$ under an applied force $F$, the gradient gives
$x_\text{eq} = F/(k_1 + k_2)$ (springs in parallel add their
stiffnesses), and the curvature $k_1 + k_2$ confirms stability. The
energy method scales to many springs and many masses with no new idea:
write the total potential energy, set its gradient to zero for the
equilibrium, and read stability off the curvature matrix.

% Figure: a mass hung between two springs (or pulled by a weight against a
% spring) and the total-potential-energy curve whose minimum locates the
% equilibrium. Left: the physical arrangement. Right: U_total(x) with its
% minimum marked.
\begin{figure}[htbp]
\centering
\begin{tikzpicture}[
  spring/.style={draw=engblack, thick},
  axis/.style={draw=engblack, -{Stealth}, thick},
  curve/.style={draw=engblue, very thick},
  partA/.style={draw=engred, very thick},
  partB/.style={draw=englime, very thick, dashed},
  guide/.style={draw=enggray, dashed, thin},
  mass/.style={draw=engblack, fill=engamber!30, minimum width=10mm, minimum height=7mm},
  label/.style={font=\small},
]
% Left panel: physical system.
% ceiling
\draw[spring] (-4.6,4) -- (-2.6,4);
\foreach \x in {-4.5,-4.2,-3.9,-3.6,-3.3,-3.0,-2.7} {
  \draw[draw=enggray, thin] (\x,4) -- (\x-0.18,4.28);
}
% spring (zig-zag) from ceiling down to the mass
\draw[spring] (-3.6,4) -- (-3.6,3.5)
  -- (-3.85,3.35) -- (-3.35,3.15) -- (-3.85,2.95) -- (-3.35,2.75)
  -- (-3.85,2.55) -- (-3.35,2.35) -- (-3.6,2.2) -- (-3.6,1.8);
% mass block
\node[mass] (m) at (-3.6,1.45) {$m$};
% weight arrow
\draw[draw=engred, -{Stealth}, very thick] (-3.6,1.0) -- (-3.6,0.3);
\node[engred, font=\scriptsize, anchor=west] at (-3.5,0.65) {$mg$};
% coordinate arrow x (downward displacement from natural length)
\draw[draw=engblue, -{Stealth}, thin] (-2.7,2.2) -- (-2.7,1.45);
\node[engblue, font=\scriptsize, anchor=west] at (-2.65,1.83) {$x$};
\node[font=\scriptsize, anchor=south] at (-3.6,4.34) {fixed support};

% Right panel: total potential energy curve.
% origin of the right plot at (0,0); axes
\draw[axis] (0,0) -- (5.4,0) node[anchor=west, label] {displacement $x$};
\draw[axis] (0,-0.3) -- (0,4.4) node[anchor=south, label] {$U(x)$};

% spring potential U_s = 1/2 k x^2 : parabola (red), scaled
% plotted on x in 0..5; U_s = 0.32*(x)^2 (in plot units)
\draw[partA] plot[domain=0:4.3, samples=120] ({\x}, {0.32*\x*\x});
\node[engred, font=\scriptsize, anchor=south west] at (3.4,3.7) {$U_\mathrm{spring}=\tfrac12 k x^2$};

% gravity potential U_g = -mg x : straight line (green dashed), scaled
% slope chosen so the well minimum sits around x=2.3
\draw[partB] plot[domain=0:5.2, samples=2] ({\x}, {-0.55*\x + 0.9});
\node[englime, font=\scriptsize, anchor=north west] at (3.9,-0.2) {$U_\mathrm{grav}=-mgx$};

% total U = 0.32 x^2 - 0.55 x + 0.9 : minimum at x = 0.55/(2*0.32) = 0.86 (plot units ~)
\draw[curve] plot[domain=0:5.0, samples=160] ({\x}, {0.32*\x*\x - 0.55*\x + 0.9});
\node[engblue, font=\scriptsize, anchor=south west] at (2.6,2.0) {$U_\mathrm{total}$};

% minimum marker: x_min = 0.859, U_min = 0.32*0.738 - 0.55*0.859 + 0.9 = 0.664
\fill[engblack] (0.859,0.664) circle (2pt);
\draw[guide] (0.859,0) -- (0.859,0.664);
\node[font=\scriptsize, anchor=north] at (0.859,-0.05) {$x_\mathrm{eq}$};
\node[font=\scriptsize, anchor=west] at (0.95,0.62) {$U'(x_\mathrm{eq})=0$};

\end{tikzpicture}
\caption[Equilibrium as the minimum of total potential energy]{A mass on
a hanging spring (left) and its total potential energy as a function of
the downward displacement $x$ (right). The spring stores
$U_\mathrm{spring}=\tfrac12 k x^2$ (red), a rising parabola; gravity
contributes $U_\mathrm{grav}=-mgx$ (green dashed), a falling straight
line. Their sum $U_\mathrm{total}$ (blue) has a single minimum at the
equilibrium displacement $x_\mathrm{eq}$, where the slope vanishes:
$U'(x_\mathrm{eq}) = k x_\mathrm{eq} - mg = 0$, hence $x_\mathrm{eq} =
mg/k$, the familiar static stretch. The principle of minimum potential
energy locates stable equilibrium at the bottom of the well and reads
stability off the curvature $U'' = k > 0$ there, recovering the
force-balance result without writing a single force equation.}
\label{fig:vol03ch05:min-pe-springs}
\end{figure}


\subsection{Worked example: work-energy accounting with friction}

Friction breaks the conservation of mechanical energy, and the
energy method handles it through a single change to the balance: the
mechanical energy at the end equals the mechanical energy at the start
minus the work dissipated by friction. A block of mass $m =
4.0\,\si{\kilogram}$ is launched up a rough incline of angle
$\alpha = 25^{\circ}$ with initial speed $v_0 = 6.0\,\si{\meter\per
\second}$; the coefficient of kinetic friction is $\mu = 0.30$. How far
up the incline does the block travel before stopping, and how fast is it
moving when it returns to the launch point?

\textbf{Going up.} The block rises a distance $d$ along the incline,
gaining height $d\sin\alpha$ and losing all its kinetic energy to gravity
and friction. The friction force is $\mu m g\cos\alpha$, opposing the
upward motion, so the energy balance from launch to the highest point is
\[
\tfrac12 m v_0^2 = m g d \sin\alpha + \mu m g \cos\alpha\, d.
\]
The mass cancels, and
\[
d = \frac{v_0^2}{2 g(\sin\alpha + \mu\cos\alpha)}
= \frac{6.0^2}{2 \cdot 9.81 (\,0.423 + 0.30\cdot 0.906\,)}
= \frac{36}{2\cdot 9.81 \cdot 0.695}
\approx 2.64\,\si{\meter}.
\]

\textbf{Coming down.} On the descent gravity drives the block down and
friction now opposes the downward motion, so the energy released by the
descent is the gravitational drop minus the friction loss over the same
distance $d$:
\[
\tfrac12 m v_f^2 = m g d \sin\alpha - \mu m g \cos\alpha\, d.
\]
Solving for the return speed,
\[
v_f = \sqrt{2 g d(\sin\alpha - \mu\cos\alpha)}
= \sqrt{2 \cdot 9.81 \cdot 2.64 (\,0.423 - 0.272\,)}
\approx 2.80\,\si{\meter\per\second}.
\]
The block returns at $2.8\,\si{\meter\per\second}$, well under its launch
speed of $6.0\,\si{\meter\per\second}$, because friction took its toll
twice, once on the way up and once on the way down. The total mechanical
energy lost to friction over the round trip is the kinetic-energy deficit
$\tfrac12 m(v_0^2 - v_f^2) = \tfrac12\cdot 4.0\,(36 - 7.84) \approx
56\,\si{\joule}$, which equals $2\mu m g\cos\alpha\,d = 2\cdot 0.30\cdot
4.0\cdot 9.81\cdot 0.906\cdot 2.64 \approx 56\,\si{\joule}$, the friction
work over the full $2d$ of travel. The two figures agree, which closes
the accounting: every joule the block started with is either returned as
kinetic energy or charged to friction.

\begin{mastery}
\textbf{Betti-Maxwell reciprocity and complementary energy.} Two
results deepen the strain-energy machinery. The first is \engterm{Betti's
reciprocal theorem}: for a linearly elastic structure under two separate
load systems $A$ and $B$, the work done by the forces of system $A$
moving through the displacements caused by system $B$ equals the work done
by the forces of system $B$ moving through the displacements caused by
system $A$. Its special case, the \engterm{Maxwell reciprocity} relation,
states that the deflection at point $i$ due to a unit load at point $j$
equals the deflection at $j$ due to a unit load at $i$, so the structure's
flexibility matrix is symmetric. The symmetry is not a convenience of
notation; it is a physical constraint that a correct stiffness or
flexibility matrix must satisfy, and a finite-element matrix that comes
out asymmetric signals a formulation error. The reciprocity follows from
the existence of a strain-energy function that is a state function of the
loads, the same property that makes Castigliano's theorem hold.

The second is the distinction between strain energy $U$, a function of
the displacements, and \engterm{complementary energy} $U^*$, a function
of the loads, related by the Legendre transform $U + U^* = \sum_i P_i
\delta_i$ exactly as the Hamiltonian and Lagrangian are related in
section 5.3. For a linearly elastic structure the force-deflection line
is straight, the two triangular areas are equal, and $U = U^*$ in
magnitude, which is why Castigliano's theorem can be stated against either
energy without confusion. For a nonlinear elastic material the two
energies differ, and the careful statement of Castigliano's first theorem
($P_i = \partial U/\partial\delta_i$, the load as the derivative of strain
energy with respect to displacement) and second theorem ($\delta_i =
\partial U^*/\partial P_i$, the displacement as the derivative of
complementary energy with respect to load) keeps the two apart. The
minimum-potential-energy principle of the spring example is the
displacement-based statement; the principle of minimum complementary
energy is its load-based dual, and the pair is the variational foundation
of the displacement and force methods of structural analysis, and of the
stiffness and flexibility formulations of the finite-element method.
\end{mastery}

\section{Failure: missing degrees of freedom in dynamic analysis}
\pathtag{core}

The energy method works only as well as the engineer's choice of
generalised coordinates. A coordinate set that fails to capture all
the system's degrees of freedom produces equations of motion that are
correct for the assumed reduced system but blind to motions of the
real system that the coordinates cannot represent. The reduced
analysis predicts the system's response to loads within the
captured modes; the real system's response includes modes the
analysis cannot see. When those uncaptured modes are excited by the
loading, the analysis fails not by computational error but by
missing the question.

The chapter's failure section names the canonical case of this
mode.

\subsection{Tacoma Narrows: a missing torsional degree of freedom}

The original Tacoma Narrows suspension bridge opened on 1 July 1940
and collapsed on 7 November 1940, four months later. On the day of
the collapse, a steady wind of approximately $19\,\si{\meter\per\
second}$ excited a torsional self-excited oscillation of the bridge's
deck; the amplitude grew until structural members failed and the
$853\,\si{\meter}$ centre span tore from its suspension cables and
fell into Puget Sound. The Federal Works Agency report of 28 March
1941, prepared by O.\ H.\ Ammann, Theodore von K\'arm\'an, and Glenn
B.\ Woodruff, identified the collapse as a case of aeroelastic
torsional flutter \cite{acc:tacoma-narrows-1941}.

The mechanism is a coupled motion of the deck and the surrounding
airflow. The deck oscillated in a torsional mode, rotating about its
longitudinal centreline; at each instantaneous angular displacement,
the deck's cross-section presented an angle of attack to the steady
wind, and the wind exerted a pitching moment on the deck that
depended on the angle of attack. At a critical wind speed, the
phase relationship between the deck's angular motion and the
aerodynamic moment became such that the wind did net positive work
on the deck over each cycle, feeding energy into the torsional
oscillation. The amplitude grew exponentially until the structural
capacity of the deck and its connections was exceeded.

The Tacoma Narrows design used the structural-dynamics methods of
its period. The deck was analysed for static wind loads (a horizontal
pressure of $\sim 2.4\,\text{kPa}$ on the projected vertical area)
and for the small-deflection vertical bending modes. The torsional
degree of freedom of the deck section in interaction with the wind
was not part of the analysis. The generalised coordinates of the
structural-dynamics model did not include the deck's torsional mode
coupled to the airflow; the analysis was, in the language of this
chapter, blind to the mode that killed the bridge.

In the analytical-mechanics framework of this chapter, the failure
is precisely a too-small generalised-coordinate set. The deck's
genuine dynamic state requires at least three generalised coordinates
per cross-section: a vertical heave, a horizontal sway, and a
torsional rotation. The aerodynamic loads on the deck depend on all
three; the coupling between deck motion and aerodynamic moment makes
the system not a passive structure under external load but a
self-excited oscillator with an aeroelastic stiffness and aeroelastic
damping. A correctly posed dynamic analysis must include the
torsional coordinate and must include the aerodynamic-feedback term
in the Lagrangian (or equivalently in the equations of motion, as a
generalised non-conservative force coupled to $\dot{\theta}_{
\text{deck}}$). The Tacoma Narrows analysis included neither.

The mechanism reduces to a single equation that the energy method makes
transparent. Model the deck's torsional coordinate $\theta$ as a damped
oscillator driven by an aerodynamic moment that, to leading order,
depends on the torsional rate:
\[
I\ddot{\theta} + c\dot{\theta} + k\theta = M_\text{aero}
\approx q\,A\,\frac{\partial C_M}{\partial\dot\theta}\,\dot\theta,
\]
with $I$ the deck's torsional inertia per unit length, $c$ the
structural damping, $k$ the torsional stiffness, and $q = \tfrac12\rho
U^2$ the dynamic pressure of a wind of speed $U$. Collecting the
rate terms gives an \emph{effective} damping coefficient
\[
c_\text{eff} = c - q\,A\,\frac{\partial C_M}{\partial\dot\theta}.
\]
As long as $c_\text{eff} > 0$ the deck damps disturbances and is stable.
The aerodynamic term grows with the dynamic pressure $q\propto U^2$, so
above a critical wind speed $U_\text{crit}$, where the aerodynamic
contribution overtakes the structural damping, $c_\text{eff}$ goes
negative. A negative damping coefficient is an oscillator that gains
energy each cycle: the amplitude grows as $e^{|c_\text{eff}|t/(2I)}$
until structural members fail. This is the self-excited flutter the
static-load analysis could not see, because the static analysis carried
no $\dot\theta$ term at all and so could not produce a $c_\text{eff}$ to
drive negative. The energy bookkeeping names the failure exactly: the
wind does net positive work on the torsional coordinate per cycle, and
the missing coordinate is the one the work was being done on.

The Federal Works Agency report did not blame Leon Moisseiff, the
bridge's designer (who had also worked on the George Washington and
Golden Gate bridges and was a respected suspension-bridge engineer
of the period); the report identified the deficiency of the
profession's design canon as the underlying mechanism. Suspension-
bridge practice in the 1930s had been moving toward shallower and
more slender deck profiles for aesthetic and economic reasons. The
aerodynamic consequences of those choices had not been brought into
routine design analysis; wind was treated as a static distributed
load. After Tacoma Narrows, suspension-bridge design incorporated
aerodynamic analysis as a standard discipline. The Tacoma Narrows
replacement bridge (1950) used a deeper truss-stiffened deck and
explicit aerodynamic analysis; subsequent suspension-bridge
practice includes wind-tunnel testing of deck section models.

\subsection{The diagnostic principle}

When an energy-method or generalised-coordinate analysis predicts a
result that is contradicted by observation, the first diagnostic
question is whether all the real system's degrees of freedom are
present in the model. The diagnostic check is to enumerate the
modes a complete free-body analysis would represent and to verify
that each appears in the generalised-coordinate set; modes not
represented are silently zeroed by the analysis. The check is
mechanical when the structure is simple (a single rigid body, a
truss) and difficult when the structure is multi-component (a
suspension bridge, a turbine-foundation system, a rocket with
propellant slosh modes). The discipline is the same: count the
degrees of freedom of the real system, then count the generalised
coordinates of the model, then identify the gap.

We return to the discipline in chapter 6 of this volume in the
context of vibration modes (the Tacoma Narrows oscillation was a
specific vibration mode with aerodynamic feedback), in chapter 13
of this volume in the context of bluff-body aerodynamics, and in
chapter 4 of Volume X as the canonical case of an unrecognised
failure mode in an established engineering practice.

\begin{principle}[Generalised coordinates must capture every mode]
An energy-method analysis with $n$ generalised coordinates can
represent at most $n$ independent modes of motion. Modes of the
real system not represented in the coordinate set are zeroed by
the analysis without warning. Before posing a Lagrangian, enumerate
the real system's degrees of freedom and verify that each is in the
coordinate set. The failure mode of an under-coordinated analysis
is not a wrong number; it is a missed phenomenon.
\end{principle}

\section{Project}\pathtag{core}

\begin{project}[Double pendulum: derive the equations of motion two
ways and reconcile]
The reader derives the equations of motion of a planar double
pendulum first by Newton's method (free-body diagrams of each bob,
explicit tensions in each string, elimination of the tensions),
then by Lagrange's method (kinetic and potential energies in the
swing-angle coordinates, Euler-Lagrange equations). The reader
verifies algebraically that the two derivations produce identical
equations, then implements a numerical simulation of the equations
and explores the motion. The deliverable is a twenty- to forty-page
analysis report.

\textbf{Track}. Derivation plus simulation. \textbf{Hazard class}:
none.

\textbf{Newtonian derivation.} The reader draws free-body diagrams of
each bob; identifies the unknown tensions $T_1$ (upper string) and
$T_2$ (lower string); writes $\vec{F} = m\vec{a}$ for each bob in
appropriate coordinates; eliminates $T_1$ and $T_2$ from the four
component equations to leave two equations of motion in $\theta_1
(t)$ and $\theta_2(t)$.

\textbf{Lagrangian derivation.} The reader writes $T$ and $U$ in the
swing-angle coordinates, forms $L = T - U$, takes the partial
derivatives, and writes out the two Euler-Lagrange equations. The
calculation follows the worked example of section 5.6.2.

\textbf{Reconciliation.} The reader puts the two pairs of equations
side by side and verifies algebraically that they are identical. The
reader writes a one-page commentary on which derivation was faster,
which was less error-prone, and which made the system's structure
(coupling, conservation laws, small-amplitude behaviour) more visible.

\textbf{Simulation.} The reader implements the equations of motion in
a numerical solver (Runge-Kutta or similar) and explores the motion
for at least four sets of initial conditions: (a) both pendula at
small amplitude in phase; (b) both at small amplitude out of phase;
(c) upper at moderate amplitude, lower at rest; (d) both at large
amplitude with energy near the over-the-top regime. For each, the
reader plots $\theta_1(t)$, $\theta_2(t)$, and the system's total
energy as a function of time (which should be constant to within
the integrator's accuracy).

\textbf{Sensitivity test.} For initial conditions (c) or (d), the
reader runs the simulation twice with initial conditions differing
by $10^{-6}$ in $\theta_1(0)$. The reader plots the difference in
$\theta_2(t)$ as a function of time and verifies that the difference
grows exponentially. The exponential growth rate is the Lyapunov
exponent of the double pendulum, a standard signature of
deterministic chaos.

\textbf{Deliverable}. The two derivations side by side, the
reconciliation, the simulation code, the four trajectory plots, the
energy-conservation plot, the sensitivity plot, and a $1{,}500$-word
reflection on what each derivation method made easy or hard to see.

\textbf{Time}. Six to twelve hours of derivation (the Newtonian
derivation is the slow step), four to eight hours of simulation
coding, two to four hours of write-up.
\end{project}

\section{Exercises}

The exercises mix calculation, derivation, estimation, simulation,
design, diagnosis, failure analysis, and judgment. Solutions are
provided for calculation and derivation exercises; estimation
exercises receive published reference estimates with the editor's
working; design, diagnosis, failure-analysis, and judgment exercises
receive rubrics rather than single answers.

\subsection*{Calculation}\pathtag{core}

\begin{exercise}
A simple pendulum of length $L = 0.5\,\si{\meter}$ swings with small
amplitude. Find its period and angular frequency.

\paragraph{Solution sketch.} The small-amplitude angular frequency is
$\omega = \sqrt{g/L} = \sqrt{9.81/0.5} \approx 4.43\,\si{\radian\per
\second}$. The period is $T = 2\pi/\omega = 2\pi\sqrt{L/g} \approx
1.42\,\si{\second}$. The result is independent of the bob mass and of
the amplitude, to the small-angle approximation.
\end{exercise}

\begin{exercise}
A bead of mass $m$ slides on a frictionless horizontal wire under no
applied force. State the bead's Lagrangian, write the Euler-Lagrange
equation, and identify the conserved quantity.

\paragraph{Solution sketch.} With $x$ the position along the wire,
$T = \tfrac12 m \dot{x}^2$ and $U = 0$, so $L = \tfrac12 m \dot{x}^2$.
Then $\partial L/\partial\dot{x} = m\dot{x}$ and $\partial L/\partial x
= 0$, so the Euler-Lagrange equation is $\tfrac{d}{dt}(m\dot{x}) = 0$,
hence $m\ddot{x} = 0$. The coordinate $x$ is cyclic, so the conjugate
momentum $p = m\dot{x}$, the linear momentum, is conserved; the bead
moves at constant velocity.
\end{exercise}

\begin{exercise}
A mass $m$ on a spring of stiffness $k$ on a frictionless horizontal
surface has the Lagrangian $L = \tfrac{1}{2} m \dot{x}^2 - \tfrac{1}{2}
k x^2$. Derive the equation of motion. Identify the natural angular
frequency.

\paragraph{Solution sketch.} $\partial L/\partial\dot{x} = m\dot{x}$ and
$\partial L/\partial x = -kx$, so the Euler-Lagrange equation
$\tfrac{d}{dt}(m\dot{x}) + kx = 0$ gives $m\ddot{x} + kx = 0$, that is
$\ddot{x} + (k/m)x = 0$. The natural angular frequency is
$\omega_0 = \sqrt{k/m}$, with period $T = 2\pi\sqrt{m/k}$.
\end{exercise}

\begin{exercise}
A pendulum of length $L$ and mass $m$ is forced by a horizontal force
$F(t) = F_0 \cos\Omega t$ applied to the bob. State the generalised
force $Q^{\text{nc}}$ in the $\theta$-coordinate equation of motion.

\paragraph{Solution sketch.} The generalised force conjugate to
$\theta$ is $Q^{\text{nc}} = \vec{F}\cdot\partial\vec{r}/\partial\theta$.
The bob's position is $\vec{r} = (L\sin\theta,\,-L\cos\theta)$, so
$\partial\vec{r}/\partial\theta = (L\cos\theta,\,L\sin\theta)$. With the
horizontal force $\vec{F} = (F_0\cos\Omega t,\,0)$, the dot product is
$Q^{\text{nc}} = F_0 L \cos\theta\,\cos\Omega t$. The equation of motion
is then $mL^2\ddot{\theta} + mgL\sin\theta = F_0 L\cos\theta\,\cos\Omega
t$; for small $\theta$ this is the forced linear oscillator
$\ddot{\theta} + (g/L)\theta \approx (F_0/mL)\cos\Omega t$.
\end{exercise}

\begin{exercise}
A block of mass $m$ slides on a frictionless inclined plane of angle
$\alpha$. Choose a generalised coordinate, write the Lagrangian,
derive the equation of motion, and find the block's acceleration.

\paragraph{Solution sketch.} Let $s$ be the distance travelled down the
incline. Then $T = \tfrac12 m \dot{s}^2$ and, with the datum at the
start, the height dropped is $s\sin\alpha$, so $U = -mgs\sin\alpha$. The
Lagrangian is $L = \tfrac12 m\dot{s}^2 + mgs\sin\alpha$. The
Euler-Lagrange equation gives $m\ddot{s} - mg\sin\alpha = 0$, hence the
acceleration $\ddot{s} = g\sin\alpha$, independent of mass and the
familiar incline result.
\end{exercise}

\begin{exercise}
A bead of mass $m$ slides on a frictionless rotating hoop of radius
$R$ rotating at $\Omega = 2\sqrt{g/R}$. Find the bead's stable
equilibrium angles on the hoop.

\paragraph{Solution sketch.} The off-axis equilibria satisfy
$\cos\theta_c = g/(R\Omega^2)$. With $\Omega^2 = 4g/R$, $R\Omega^2 = 4g$,
so $\cos\theta_c = g/(4g) = 1/4$, giving $\theta_c \approx \pm 75.5^\circ$.
Since $R\Omega^2 = 4g > g$, the bottom $\theta = 0$ is unstable and the
two off-axis points $\pm 75.5^\circ$ are the stable equilibria. (The
top $\theta = \pi$ is always unstable.) This matches the
$\Omega/\Omega_\text{crit} = 2$ row of \path{data/hoop_equilibria.csv}.
\end{exercise}

\begin{exercise}
A double pendulum has equal masses $m$ and equal string lengths $L$.
Write the small-amplitude (linearised) equations of motion. Find the
two natural frequencies of small-amplitude oscillation.

\paragraph{Solution sketch.} Linearising the equal-mass equal-length
double pendulum gives $2\ddot{\theta}_1 + \ddot{\theta}_2 +
2(g/L)\theta_1 = 0$ and $\ddot{\theta}_1 + \ddot{\theta}_2 +
(g/L)\theta_2 = 0$. Substituting $\theta_i = a_i\cos\omega t$ yields the
eigenvalue problem $(K - \omega^2 M)a = 0$ with $M =
\bigl(\begin{smallmatrix}2&1\\1&1\end{smallmatrix}\bigr)$ and
$K = (g/L)\bigl(\begin{smallmatrix}2&0\\0&1\end{smallmatrix}\bigr)$. The
determinant condition gives $\omega^2 = (2\mp\sqrt2)\,g/L$, so
$\omega_1 = \sqrt{(2-\sqrt2)g/L} \approx 0.765\sqrt{g/L}$ (in-phase) and
$\omega_2 = \sqrt{(2+\sqrt2)g/L} \approx 1.848\sqrt{g/L}$ (out-of-phase),
matching \path{data/double_pendulum_modes.csv}.
\end{exercise}

\subsection*{Derivation}\pathtag{standard}

\begin{exercise}
Derive the Euler-Lagrange equation from Hamilton's principle.
State the boundary-condition assumption and identify where it is
used in the integration-by-parts step.

\paragraph{Solution sketch.} Vary the action $S = \int_{t_1}^{t_2}
L(q,\dot q,t)\,dt$ by $q\to q+\delta q$ with $\delta q(t_1) = \delta
q(t_2) = 0$ (the fixed-endpoint assumption). To first order, $\delta S =
\int [\,(\partial L/\partial q)\delta q + (\partial L/\partial\dot
q)\delta\dot q\,]\,dt$. Integrate the second term by parts:
$\int (\partial L/\partial\dot q)\,\delta\dot q\,dt = [(\partial
L/\partial\dot q)\delta q]_{t_1}^{t_2} - \int \tfrac{d}{dt}(\partial
L/\partial\dot q)\,\delta q\,dt$. The boundary term vanishes precisely
because $\delta q$ is zero at both endpoints; this is where the
boundary-condition assumption is used. Then $\delta S = \int [\,\partial
L/\partial q - \tfrac{d}{dt}(\partial L/\partial\dot q)\,]\delta q\,dt$,
and requiring $\delta S = 0$ for arbitrary $\delta q$ forces the
bracket to vanish, giving $\tfrac{d}{dt}(\partial L/\partial\dot q) -
\partial L/\partial q = 0$.
\end{exercise}

\begin{exercise}
Show that the Euler-Lagrange equation for a Lagrangian $L = \tfrac{1}{2}
m \dot{q}^2 - U(q)$ in a single generalised coordinate $q$ reduces to
Newton's second law for a particle moving in the potential $U$.

\paragraph{Solution sketch.} Here $\partial L/\partial\dot q = m\dot q$,
so $\tfrac{d}{dt}(\partial L/\partial\dot q) = m\ddot q$, and $\partial
L/\partial q = -dU/dq$. The Euler-Lagrange equation $\tfrac{d}{dt}
(\partial L/\partial\dot q) - \partial L/\partial q = 0$ becomes $m\ddot
q + dU/dq = 0$, that is $m\ddot q = -dU/dq = F(q)$. This is Newton's
second law for a particle of mass $m$ in the potential $U$, with the
conservative force $F = -dU/dq$.
\end{exercise}

\begin{exercise}
A particle moves in a central force field with potential $U(r)$.
Choose polar coordinates $(r, \phi)$. Write the Lagrangian. Identify
the cyclic coordinate and the corresponding conserved momentum. State
the physical meaning of the conserved momentum.

\paragraph{Solution sketch.} In polar coordinates the velocity squared
is $\dot r^2 + r^2\dot\phi^2$, so $T = \tfrac12 m(\dot r^2 +
r^2\dot\phi^2)$ and $L = \tfrac12 m(\dot r^2 + r^2\dot\phi^2) - U(r)$.
The coordinate $\phi$ does not appear in $L$ (only $\dot\phi$ does), so
$\phi$ is cyclic and $p_\phi = \partial L/\partial\dot\phi =
mr^2\dot\phi$ is conserved. This is the angular momentum about the
force centre; its conservation is the rotational symmetry of the central
potential, equivalent to Kepler's law of equal areas in equal times.
\end{exercise}

\begin{exercise}
Derive the small-amplitude (linearised) equations of motion of a
double pendulum with equal masses and equal string lengths. Identify
the two normal modes and their natural frequencies.

\paragraph{Solution sketch.} Linearising as in the calculation-block
companion gives $M = \bigl(\begin{smallmatrix}2&1\\1&1\end{smallmatrix}
\bigr)$, $K = (g/L)\bigl(\begin{smallmatrix}2&0\\0&1\end{smallmatrix}
\bigr)$. The eigenvectors of $M^{-1}K$ are the mode shapes. The
low-frequency mode $\omega_1^2 = (2-\sqrt2)g/L$ has the bobs in phase
with amplitude ratio $a_2/a_1 = +\sqrt2/2$; the high-frequency mode
$\omega_2^2 = (2+\sqrt2)g/L$ has them out of phase with $a_2/a_1 =
-\sqrt2$. The lower link swings with larger amplitude than the upper in
both modes. General small-amplitude motion is a superposition of the
two, each oscillating at its own frequency; the resulting beat pattern
is what makes a lightly disturbed double pendulum appear to trade energy
between the links.
\end{exercise}

\begin{exercise}
Derive the bifurcation condition $R\Omega^2 = g$ for the bead-on-
rotating-hoop system of section 5.6.3. State the physical
interpretation of the bifurcation.

\paragraph{Solution sketch.} The equilibrium condition is
$\sin\theta\,(g - R\Omega^2\cos\theta) = 0$. The off-axis solutions
$\cos\theta_c = g/(R\Omega^2)$ require the right side to lie in
$(0,1)$, which is possible only when $R\Omega^2 > g$. As $\Omega$ rises
through the critical value $\Omega_\text{crit} = \sqrt{g/R}$, the
condition $R\Omega^2 = g$ marks the appearance of the two new
equilibria, simultaneously with the bottom losing stability. The
interpretation: below the threshold the gravitational restoring effect
dominates and the bead rests at the bottom; above it the centrifugal
effect of the rotation exceeds gravity at the bottom, the bottom becomes
unstable, and the bead is flung outward to a new stable angle. The
transition is a pitchfork bifurcation, a symmetry-breaking event in
which the single symmetric equilibrium gives way to a symmetric pair.
\end{exercise}

\subsection*{Estimation}\pathtag{core}

\begin{exercise}
Estimate the period of a $1\,\si{\meter}$ pendulum swinging at
amplitude $\theta_0 = 45^{\circ}$. State the correction to the small-
amplitude formula and bracket the answer.

\paragraph{Reference estimate.} Small-angle period $T_0 = 2\pi\sqrt{1/
9.81} \approx 2.006\,\si{\second}$. With $\theta_0 = 45^\circ = 0.785\,
\si{\radian}$, the first correction $T \approx T_0(1 + \theta_0^2/16) =
2.006\,(1 + 0.0386) \approx 2.083\,\si{\second}$. The exact value from
\path{data/pendulum_period.csv} is $2.086\,\si{\second}$, so the true
period sits between the small-angle $2.01\,\si{\second}$ and roughly
$2.09\,\si{\second}$, about $4\,\%$ above the small-angle figure. Bracket:
$2.01$ to $2.09\,\si{\second}$.
\end{exercise}

\begin{exercise}
Estimate the energy stored in the elastic deformation of a typical
diving board when supporting a $70\,\si{\kilogram}$ diver at full
deflection. State assumptions about board geometry and material.

\paragraph{Reference estimate.} Treat the board as a cantilever
deflecting by $\delta \approx 0.15\,\si{\meter}$ under the diver's
weight $W = mg \approx 686\,\si{\newton}$ at the tip. A linear-spring
model stores $U = \tfrac12 W\delta$ at full static deflection (the
factor of one-half because the force builds up linearly with the
deflection), so $U \approx \tfrac12 \cdot 686 \cdot 0.15 \approx
51\,\si{\joule}$. Order of magnitude: tens of joules, comparable to the
$mg\,h$ a diver gains by rising $\sim 0.1\,\si{\meter}$. Assumptions: a
springboard tip deflection of order $0.1$ to $0.2\,\si{\meter}$ and a
roughly linear load-deflection relation up to full deflection. A
fibreglass board returning most of this energy is what gives the
spring its launch.
\end{exercise}

\begin{exercise}
Estimate the critical rotation speed at which a bead on a $30\,\si
{\milli\meter}$-radius rotating hoop transitions from the bottom-
equilibrium regime to the off-axis-equilibrium regime. Express in
revolutions per minute.

\paragraph{Reference estimate.} The threshold is $R\Omega^2 = g$, so
$\Omega_\text{crit} = \sqrt{g/R} = \sqrt{9.81/0.030} \approx
18.1\,\si{\radian\per\second}$. Converting, $f = \Omega/(2\pi) \approx
2.88$ revolutions per second, or about $173$ revolutions per minute.
This matches the value printed by \path{code/hoop_bifurcation.py}.
Below roughly $170$ rev/min the bead sits at the bottom; above it the
bead climbs to an off-axis equilibrium.
\end{exercise}

\subsection*{Simulation}\pathtag{mastery}

\begin{exercise}
Implement a numerical Runge-Kutta solver for the simple-pendulum
equation $\ddot{\theta} + (g/L)\sin\theta = 0$. Compare the simulated
period to the small-amplitude prediction for swing amplitudes
$10^{\circ}$, $30^{\circ}$, $60^{\circ}$, and $90^{\circ}$. Plot the
period as a function of amplitude.

\paragraph{Solution sketch.} The reference implementation is
\path{code/pendulum_period.py}: integrate $\ddot\theta + (g/L)\sin\theta
= 0$ as a first-order system from rest at $\theta_0$, and measure the
period from the time to return to the turning point (or four times the
time to the first downward zero crossing). Expected results, from
\path{data/pendulum_period.csv} with $L = 1\,\si{\meter}$, $T_0 \approx
2.006\,\si{\second}$: at $10^\circ$ the period exceeds $T_0$ by about
$0.19\,\%$; at $30^\circ$ by $1.7\,\%$; at $60^\circ$ by $7.3\,\%$; at
$90^\circ$ by $18\,\%$. The period rises monotonically with amplitude
and diverges logarithmically as $\theta_0\to 180^\circ$. The first
correction $1 + \theta_0^2/16$ tracks the curve to better than $0.5\,\%$
out to about $60^\circ$.
\end{exercise}

\begin{exercise}
Implement a Lagrangian-mechanics solver for the double pendulum.
Verify energy conservation to within $0.1\,\%$ over a $30\,\si{\second}$
simulation. Compare the trajectories for nearby initial conditions
and demonstrate exponential divergence (the chapter's project takes
this further).

\paragraph{Solution sketch.} The reference implementation is
\path{code/double_pendulum.py}: form the two-by-two mass matrix from the
linearised-acceleration coefficients, invert it at each step to extract
$(\ddot\theta_1,\ddot\theta_2)$, advance with RK4, and log the total
energy $E = T + U$ at each step. A correctly implemented conservative
integrator with a step of $10^{-3}\,\si{\second}$ holds $E$ constant to
under $2\times 10^{-7}$ relative over $30\,\si{\second}$, comfortably
inside the $0.1\,\%$ target, as recorded in
\path{data/double_pendulum_trajectory.csv}. Repeating a large-amplitude
run with $\theta_1(0)$ perturbed by $10^{-6}\,\si{\radian}$ produces a
separation $|\Delta\theta_2(t)|$ that grows exponentially, a straight
line on a log scale (figure~\ref{fig:vol03ch05:double-pendulum-divergence}),
until it saturates near one radian. Energy drift growing faster than
this bound signals a coding error in the mass-matrix inversion or too
large a step.
\end{exercise}

\subsection*{Design}\pathtag{standard}

\begin{exercise}
Design a pendulum clock with period exactly $2\,\si{\second}$.
Specify the pendulum length and the regulating-mechanism principle.
State the temperature-induced length change that would shift the
clock by $1\,\si{\second\per\day}$.

\paragraph{Model-answer sketch.} For $T = 2\,\si{\second}$, the
small-amplitude relation $L = g(T/2\pi)^2$ gives $L = 9.81\cdot
(1/\pi)^2 \approx 0.994\,\si{\meter}$, the classic ``seconds pendulum''
of just under a metre (its half-period, one swing, is one second). The
regulating mechanism is an escapement that delivers a small impulse each
swing to replace damping losses, keeping the amplitude small and
constant so the isochronism of the small-angle period holds. Sensitivity:
$T\propto\sqrt L$, so $\Delta T/T = \tfrac12\,\Delta L/L$. A drift of
$1\,\si{\second}$ per day is $\Delta T/T = 1/86400 \approx
1.16\times 10^{-5}$, hence $\Delta L/L = 2.3\times 10^{-5}$, or $\Delta
L \approx 23\,\si{\micro\meter}$ on a one-metre rod. A steel rod
($\alpha \approx 12\times 10^{-6}\,\si{\per\kelvin}$) reaches this with
a temperature change of about $2\,\si{\kelvin}$, which is why precision
clocks use temperature-compensated gridiron or invar pendulums.
\end{exercise}

\begin{exercise}
Design a torsion pendulum (rotor on a thin wire) to measure small
torques. Specify the rotor's moment of inertia, the wire's torsion
stiffness, and the natural period of the pendulum. State the minimum
torque the pendulum can detect, given a $1$-arcsecond angular
resolution.

\paragraph{Model-answer sketch.} The torsion pendulum obeys $I\ddot
\phi + \kappa\phi = \tau_\text{ext}$, with $I$ the rotor moment of
inertia and $\kappa$ the wire torsional stiffness; the natural period is
$T = 2\pi\sqrt{I/\kappa}$. A worked choice: a rotor of $I \approx
10^{-4}\,\si{\kilogram\meter\squared}$ on a wire of $\kappa \approx
10^{-6}\,\si{\newton\meter\per\radian}$ gives $T = 2\pi\sqrt{10^{-4}/
10^{-6}} = 2\pi\cdot 10 \approx 63\,\si{\second}$, a slow, easily read
oscillation. The detectable torque follows from the static deflection
$\tau = \kappa\,\phi_\text{min}$: with $\phi_\text{min} =
1\,\text{arcsec} = 4.85\times 10^{-6}\,\si{\radian}$, the minimum torque
is $\tau \approx 10^{-6}\cdot 4.85\times 10^{-6} \approx 5\times
10^{-12}\,\si{\newton\meter}$. The design trade is explicit: a softer
wire (smaller $\kappa$) detects smaller torques but lengthens the period
and raises sensitivity to drift and seismic noise. This is the
Cavendish-balance principle.
\end{exercise}

\begin{exercise}
Design a vibration test rig for evaluating the natural frequency of
a small cantilever beam. State the generalised coordinates required,
the energy expressions, and the predicted natural frequency for a
$300\,\si{\milli\meter}$ aluminum beam of $1\,\si{\milli\meter}$
thickness.

\paragraph{Model-answer sketch.} One generalised coordinate suffices
for the fundamental mode: the tip deflection $w(t)$, with an assumed
mode shape $\phi(x)$ satisfying the clamped-free boundary conditions.
Kinetic energy $T = \tfrac12 m_\text{eff}\dot w^2$ and elastic potential
$U = \tfrac12 k_\text{eff} w^2$ give $\omega_1 = \sqrt{k_\text{eff}/
m_\text{eff}}$. The exact Euler-Bernoulli result is $\omega_1 =
(1.875)^2\sqrt{EI/(\rho A L^4)}$. For an aluminium beam ($E \approx
69\,\si{\giga\pascal}$, $\rho \approx 2700\,\si{\kilogram\per\meter
\cubed}$) of width $b$, thickness $h = 1\,\si{\milli\meter}$, length
$L = 0.3\,\si{\meter}$, the width cancels: $I/A = h^2/12$, so $\omega_1 =
3.516\sqrt{E h^2/(12\rho L^4)} = 3.516\sqrt{69\times 10^9\cdot 10^{-6}/
(12\cdot 2700\cdot 0.3^4)} \approx 3.516\cdot 8.2 \approx
29\,\si{\radian\per\second}$, that is $f_1 \approx 4.6\,\si{\hertz}$. The
rig drives the beam base near this frequency and detects the resonant
amplitude peak. A Rayleigh estimate with a parabolic trial shape lands
within a few per cent of this.
\end{exercise}

\begin{exercise}
Design a Foucault pendulum installation in a building atrium to
demonstrate Earth's rotation. Specify the pendulum length, the
swing amplitude, the expected precession rate at the installation's
latitude, and the indicator-pin spacing on the demonstration
fixture.

\paragraph{Model-answer sketch.} A long pendulum keeps the swing plane
visibly fixed while the building rotates beneath it. Choose $L \approx
20\,\si{\meter}$ for a period $T = 2\pi\sqrt{20/9.81} \approx
9\,\si{\second}$ and a heavy bob to fight air damping; swing amplitude a
few degrees. The precession rate of the swing plane is $\Omega_p =
\Omega_\oplus\sin\lambda$, with $\Omega_\oplus = 2\pi$ per sidereal day
$= 15.04^\circ$ per hour and $\lambda$ the latitude. At $\lambda =
45^\circ$, $\Omega_p \approx 10.6^\circ$ per hour, a full turn in about
$34\,\si{\hour}$. Over one swing of amplitude $A \approx 2\,\si{\meter}$
at the floor, the plane advances by $A\,\Omega_p\,T \approx 2\cdot
(10.6\cdot\pi/180/3600)\cdot 9 \approx 9\times 10^{-4}\,\si{\meter}$, so
indicator pins spaced a few millimetres apart are knocked over one per
many minutes, a readable demonstration. State the installation latitude
explicitly, since the rate vanishes at the equator and is maximal at the
poles.
\end{exercise}

\subsection*{Diagnosis and reverse engineering}\pathtag{core}

\begin{exercise}
A flexible-joint robot arm is analysed using rigid-link approximations
and exhibits unexpected oscillations at the end-effector under fast
moves. Identify the missing degree of freedom and state the modeling
correction.

\paragraph{Solution sketch.} The rigid-link model assigns one
generalised coordinate per joint and treats each link as
infinitely stiff. The missing degree of freedom is the elastic
deflection of the joint (and of the link), a compliance the rigid model
zeroes. Under slow moves the joint stiffness keeps the deflection
negligible; under fast moves the inertial torque excites the
joint-compliance mode, and the end-effector rings at the natural
frequency $\sqrt{k_\text{joint}/I_\text{link}}$. The correction is to add
a coordinate for each flexible joint, modelling it as a torsional spring
$k_\text{joint}$ in series with the actuator, giving a flexible-joint
Lagrangian with twice as many coordinates as the rigid model. The
controller must then either be bandwidth-limited below the compliance
mode or include explicit vibration damping.
\end{exercise}

\begin{exercise}
A car's suspension is analysed as a two-degree-of-freedom system
(sprung mass and unsprung mass on a tire spring). The car exhibits an
unexplained vibration mode in the $5\,\si{\hertz}$ range. Identify
the most probable missing degree of freedom and state the
diagnostic measurement.

\paragraph{Solution sketch.} The two-degree-of-freedom quarter-car model
captures the sprung-mass (body) bounce near $1\,\si{\hertz}$ and the
unsprung-mass (wheel) hop near $10$ to $15\,\si{\hertz}$. A mode near
$5\,\si{\hertz}$ falls between these and is most probably a structural
flexible mode the lumped model omits: body torsional or bending
flexibility, a subframe resonance, or seat-and-occupant dynamics. The
missing degree of freedom is the chassis or subframe elastic mode, a
coordinate the rigid-body two-mass model cannot represent. Diagnostic
measurement: an instrumented modal test (accelerometers along the body
with a shaker or impact hammer) to extract the mode shape at
$5\,\si{\hertz}$; if the shape is a chassis bending or subframe motion
rather than rigid body-on-suspension motion, the lumped model needs the
extra coordinate.
\end{exercise}

\begin{exercise}
A rocket's attitude control uses a model that treats the propellant
as a rigid mass at the centre of the tank. The control loop exhibits
limit cycles when the tank is partially full. Identify the missing
degree of freedom and state the engineering correction.

\paragraph{Solution sketch.} Treating the propellant as a rigid mass
zeroes the sloshing of the free liquid surface. The missing degree of
freedom is the propellant slosh mode: in a partially full tank the
liquid surface can rock, and to first order this behaves like a pendulum
(or spring-mass) of natural frequency set by the tank radius and fill
level. The slosh mass couples to the vehicle's attitude through the
tank-wall reaction, and when the slosh frequency lies near the control
loop's bandwidth the feedback sustains a limit cycle. The correction is
to add a slosh coordinate, modelled as an equivalent pendulum or
spring-mass attached to the vehicle, into the attitude-dynamics
Lagrangian, and to add baffles inside the tank that raise the slosh
damping and shift its frequency away from the control bandwidth. This is
the classic propellant-slosh problem of launch-vehicle guidance.
\end{exercise}

\subsection*{Failure analysis}\pathtag{core}

\begin{exercise}
Re-read the Federal Works Agency report on the Tacoma Narrows Bridge
collapse \cite{acc:tacoma-narrows-1941}. Identify the generalised
coordinates that a corrected aeroelastic analysis would include and
state which aerodynamic feedback term in the Lagrangian (or equivalent
generalised force) drove the instability.

\paragraph{Solution sketch.} A corrected aeroelastic model needs, per
deck cross-section, at least three generalised coordinates: vertical
heave $h$, horizontal sway $s$, and torsional rotation $\alpha$ about the
deck's longitudinal axis. The original analysis carried only the
vertical bending coordinate under static wind load and omitted $\alpha$.
The destabilising term is the aerodynamic pitching moment that depends on
the torsional rate $\dot\alpha$, entering as a generalised
non-conservative force $Q_\alpha = \partial M_\text{aero}/\partial\dot
\alpha\cdot\dot\alpha$ whose sign over a cycle is such that the wind does
net positive work on the deck. This negative aerodynamic damping (an
effective damping coefficient driven below zero) makes the torsional
oscillation self-excited: the amplitude grows exponentially rather than
decays. The report's diagnosis of aeroelastic torsional flutter is
exactly the statement that the $\alpha$ coordinate and its
aerodynamic-feedback term were absent from the design model
\cite{acc:tacoma-narrows-1941}.
\end{exercise}

\begin{exercise}
A wind-loaded tall structure (chimney, antenna mast, slender tower)
exhibits unexpected vortex-induced oscillation at moderate wind
speeds. State, in $200$ to $300$ words, the generalised coordinates
that a complete analysis must include and the aerodynamic-coupling
term that drives the response.

\paragraph{Model-answer sketch.} A complete analysis carries the
cross-wind transverse displacement of the structure (and, for tall
towers, the first few bending-mode coordinates) rather than treating
wind as a static drag load. The driving mechanism is vortex shedding:
the alternating wake vortices apply a transverse force at the Strouhal
frequency $f_s = St\,U/D$, with $St \approx 0.2$ for a circular section,
$U$ the wind speed, $D$ the diameter. When $f_s$ approaches a structural
natural frequency, the oscillation and the shedding lock in: the
structure's own motion organises the wake, so the forcing term becomes a
function of the response (a motion-dependent generalised force), and the
amplitude can grow until limited by nonlinearity. The coordinate that
matters is the transverse modal displacement; the coupling term is the
lift force synchronised to that displacement during lock-in. A complete
model adds aerodynamic damping that can go negative over the lock-in
range. A strong answer names lock-in explicitly and distinguishes it from
ordinary forced resonance.
\end{exercise}

\begin{exercise}
A bridge cable-stay system, designed for static wind loads only,
exhibits significant cable galloping in icing conditions. State the
missing degree of freedom in the static-load analysis and the
mechanism by which icing changes the aerodynamic-coupling term.

\paragraph{Solution sketch.} The static-load analysis omits the cable's
transverse oscillation coordinate (and, where relevant, its torsional
rotation), treating wind as a fixed drag. Galloping is a one-degree-of-
freedom aeroelastic instability of the transverse coordinate. A bare
cable is nearly circular and aerodynamically stable; an ice accretion
gives the cross-section a non-circular profile (a D-shape or crescent)
whose lift coefficient varies steeply with angle of attack. The Den
Hartog galloping criterion is that instability occurs when $dC_L/d\alpha
+ C_D < 0$, that is when the slope of the lift curve plus the drag turns
the effective aerodynamic damping of the transverse coordinate negative.
Icing changes the coupling term by reshaping the section so that
$dC_L/d\alpha$ becomes large and negative; the wind then feeds energy
into the transverse mode and the cable gallops at large, low-frequency
amplitude. The correction is aerodynamic (de-icing, section shaping) or
structural (added damping, cross-ties between cables).
\end{exercise}

\subsection*{Judgment}\pathtag{standard}

\begin{exercise}
A junior engineer asks whether Lagrangian mechanics replaces Newtonian
mechanics or supplements it. Write a one-paragraph reply that
addresses when each method is preferred and when they are
interchangeable.

\paragraph{Model-answer sketch.} Lagrangian mechanics supplements
rather than replaces Newtonian mechanics: the two are mathematically
equivalent, derivable from each other through d'Alembert's principle, so
they always give the same equations of motion. The choice is tactical.
Newton's method is preferred when the constraint and contact forces are
themselves wanted (bearing loads, joint reactions, tensions), since
those forces are eliminated in the Lagrangian formulation and must be
recovered separately. The Lagrangian method is preferred for systems
with many degrees of freedom, complicated holonomic constraints, or
non-Cartesian coordinates, where free-body bookkeeping is error-prone and
the scalar-energy route is cleaner. For a single particle or a simple
rigid body the two are interchangeable and the engineer picks whichever
is faster. A strong reply notes that the methods are one physics in two
notations.
\end{exercise}

\begin{exercise}
A senior colleague argues that energy methods are an academic
luxury and that working engineers should use free-body diagrams
exclusively. Write a one-paragraph reply that takes the argument
seriously and states the design situations in which energy methods
are essential.

\paragraph{Model-answer sketch.} The colleague is right that for many
everyday problems, a single rigid body, a statically determinate
structure, a simple machine element, a free-body diagram is the fastest
and clearest tool, and the senior engineer's fluency with it is
valuable. The argument fails as a general policy. Energy methods become
essential when the system has many degrees of freedom and many internal
constraints (multi-link robots, vehicle multibody dynamics), where
free-body bookkeeping grows quadratically and the unknown internal
forces swamp the calculation; when the analysis is variational by nature
(finite-element structural dynamics, Rayleigh-Ritz natural-frequency
estimates, optimal-control trajectory problems); and when the failure
mode is a missing degree of freedom, where the discipline of writing one
scalar energy and enumerating coordinates is the safeguard the
Tacoma Narrows designers lacked. A strong reply concedes the practical
point and bounds it: free-body diagrams for the simple and the
local, energy methods for the many-coordinate and the variational.
\end{exercise}

\begin{exercise}
Identify a mechanical system in your immediate environment whose
dynamics you have never analysed in the Lagrangian formulation
(spring-loaded door, oscillating fan, vibrating washing machine).
State, in $200$ to $300$ words, the generalised coordinates you
would choose, the kinetic and potential energy expressions, and
whether the system has any cyclic coordinates that imply a
conservation law.

\paragraph{Model-answer sketch.} The answer is open, but a strong one
follows the chapter's procedure for whatever system is chosen. Take a
spring-loaded screen door: one generalised coordinate, the opening angle
$\theta$; kinetic energy $T = \tfrac12 I\dot\theta^2$ with $I$ the
door's moment of inertia about the hinge; potential energy combining the
closer-spring term $\tfrac12 k_\theta\theta^2$ and, if the door is not
vertical-axis, a gravity term. The damper makes the system
non-conservative, so the generalised damping force enters as
$Q = -c\dot\theta$. There is no cyclic coordinate here, since $\theta$
appears in $U$, so no conservation law; the energy decays through the
damper. A reply that instead picks a system with a genuine symmetry (a
freely spinning fan coasting down, where the spin angle is cyclic and
angular momentum is conserved until friction acts) should identify the
cyclic coordinate and the conserved momentum explicitly. The graded
skill is correct identification of coordinates, energies, dissipation,
and any symmetry, not the particular object.
\end{exercise}
